% SP1.N4. Las cifras proceden de las campañas N4 v1--v4.

\newcommand{\NFourVThreeWorlds}{1\,200}
\newcommand{\NFourFiiMismatchPct}{74{,}8}
\newcommand{\NFourFiiRankRho}{0{,}32}
\newcommand{\NFourPdConvergencePct}{91{,}9}
\newcommand{\NFourPdKktMedian}{4{,}97 \times 10^{-4}}
\newcommand{\NFourPdGapPct}{9{,}04}
\newcommand{\NFourSmithGapPct}{18{,}57}
\newcommand{\NFourDynamicWorlds}{30}
\newcommand{\NFourWarmCost}{2{,}590}
\newcommand{\NFourMyopicCost}{2{,}823}
\newcommand{\NFourOracleCost}{2{,}180}
\newcommand{\NFourPdRiskDiffPct}{5{,}50}
\newcommand{\NFourOracleWarmDiff}{-0{,}415}

\newcommand{\NFourHStarWorlds}{1\,200}
\newcommand{\NFourHStarRatioWorlds}{1\,200}
\newcommand{\NFourHStarTwoPct}{96{,}3}
\newcommand{\NFourHStarThreePct}{0{,}3}
\newcommand{\NFourHStarGtThreePct}{3{,}3}
\newcommand{\NFourHStarHardLowCvPct}{13}
\newcommand{\NFourHStarHardHighCvPct}{1{,}3}
\newcommand{\NFourHStarSearchMedianMs}{9{,}2}
\newcommand{\NFourHStarPairGainPct}{34{,}5}
\newcommand{\NFourHStarTripleGainPct}{42{,}3}
\newcommand{\NFourHStarGtThreeLowRatioPct}{11{,}5}
\newcommand{\NFourHStarGtThreeHighRatioPct}{0{,}2}


% --------------------------------------------------------------------------
% Página 27 - problema común y taxonomía
% --------------------------------------------------------------------------
\leveltag{RESULTADOS SP1 · NIVEL 4 DE 4 · PROBLEMA Y TAXONOMÍA}
\section*{SP1.N4: localidad estratégica en el reclutamiento atómico}
\enlargethispage{5\baselineskip}

N4 comienza en el punto donde termina N3: una asignación entera obtenida con
información vecinal puede ser factible y, aun así, quedar atrapada en un mínimo
local. N3 estudió la información disponible; aquí se pregunta cuántos robots
deben cambiar juntos para mejorar la decisión.

La salida física no cambia. Cada robot debe terminar inactivo o asociado a una
sola carga, y toda carga activa $k\in\mathcal K_{\mathrm{act}}$ debe reunir
capacidad suficiente. En las campañas estáticas,
$\mathcal K_{\mathrm{act}}=\{1,\ldots,K\}$:
\begin{equation}
\label{eq:n4-common-problem}
\begin{gathered}
\begin{aligned}
a_i&\in\{0,1,\ldots,K\}, &
Q_k(a)&=\sum_i c_i^{\mathrm{pay}}\mathbf 1\{a_i=k\},\\[-1mm]
D_4(a)&=\sum_{k\in\mathcal K_{\mathrm{act}}}[m_k-Q_k(a)]_+, &
J_4(a)&=\sum_{i:a_i\ne0}d_{i,a_i}.
\end{aligned}
\\[-1mm]
D_4(a)=0,\qquad J_4(a)\rightarrow\min .
\end{gathered}
\end{equation}

El subíndice identifica los funcionales de N4 sobre asignaciones enteras $a$.
La propuesta principal es el juego potencial geométrico de cuotas (Geo-QPG,
\emph{geometric quota-potential game}); QPG designa su juego base.

F-I, el juego atómico Geo-QPG, constituye la contribución principal. El
resultado central del nivel reúne la identidad de potencial, el refinamiento
BR--2BR--C3 y la medición del primer orden que permite escapar. F-II y F-III son
comparadores continuos a los que se aplica el mismo cierre entero; F-IV estudia
por separado la decisión temporal.
De los 1\,200 terminales BR examinados en E9, 1\,156 admitieron primero una
mejora bilateral. El dato identifica el orden del primer escape, no cuánto
beneficio producen los pares ni qué ocurre por encima de tres robots. La
Figura~\ref{fig:n4-family-taxonomy} ordena estas funciones y sus límites.

\begin{figure}[H]
\centering
\scalebox{.70}{\pgfdeclarelayer{n4taxonomyback}
\pgfsetlayers{n4taxonomyback,main}
\begin{tikzpicture}[x=1cm,y=1cm,>=Latex,
 % anchor=north: las cuatro familias tienen alturas de contenido distintas; al
 % centrarlas en y quedaban escalonadas y su borde superior invadía la barra.
 family/.style={draw=SPGrid,rounded corners=1.1mm,text width=2.42cm,
  minimum height=2.75cm,align=center,inner sep=1.6mm,font=\scriptsize,
  anchor=north,
  % Cajas estrechas: sin esto TeX parte POBLACIONAL, RESTRINGIDA o coalición.
  execute at begin node={\hyphenpenalty=10000\exhyphenpenalty=10000\relax}},
 bar/.style={draw=SPGrid,rounded corners=1mm,fill=SPFloor,text width=10.65cm,
  minimum height=.95cm,align=center,font=\footnotesize},
 arr/.style={->,thick,draw=VIUOrange,shorten >=1.0mm}]
\node[bar,draw=SPBlue!65,fill=SPBlue!4] (problem) at (0,4.05)
 {\textbf{Mismo problema físico}\\[-.4mm]
  cuotas heterogéneas: $D(a)=0$; después minimizar $J(a)$};

\node[family,draw=SPBlue!70,fill=SPBlue!4] (f1) at (-4.35,2.62)
 {\textbf{F-I · ATÓMICA}\\[-1mm]
  \rule{2.15cm}{.2pt}\\
  \textbf{variable}\ $a_i$ entero\\
  \textbf{evolución}\ mejoras $h$-locales\\
  \textbf{objetivo}\ mínimo local de $\Psi_4$\\
  \textbf{salida}\ coalición directa\\
  \textcolor{SPGreen}{E4--E6/E9:\\[-.4mm]validada}};
\node[family,draw=SPGreen!70,fill=SPGreen!4] (f2) at (-1.45,2.62)
 {\textbf{F-II · POBLACIONAL}\\[-1mm]
  \rule{2.15cm}{.2pt}\\
  \textbf{variable}\ $x_i\in\Delta_i$\\
  \textbf{evolución}\ Replicator/Smith\\BNN/Logit\\
  \textbf{objetivo}\ estado estacionario\\
  \textbf{salida}\ $\mathcal R(x)$\\
  \textcolor{SPGreen}{E7: contrastada}};
\node[family,draw=SPPurple!70,fill=SPPurple!4] (f3) at (1.45,2.62)
 {\textbf{F-III · RESTRINGIDA}\\[-1mm]
  \rule{2.15cm}{.2pt}\\
  \textbf{variable}\ $(x,\lambda)$\\
  \textbf{evolución}\ primal--dual + consenso\\
  \textbf{objetivo}\ vGNE\\
  \textbf{salida}\ $\mathcal R(x)$\\
  \textcolor{SPGreen}{E7: contrastada}};
\node[family,draw=VIUOrange!80,fill=VIUOrange!4] (f4) at (4.35,2.62)
 {\textbf{F-IV · TEMPORAL}\\[-1mm]
  \rule{2.15cm}{.2pt}\\
  \textbf{variable}\ $(s_t,u_t)$\\
  \textbf{evolución}\ transición + política\\
  \textbf{objetivo}\ coste finito recurrente\\
  \textbf{salida}\ coalición recurrente\\
  \textcolor{SPGreen}{E8: ESTUDIO TEMPORAL}};

\node[bar,draw=SPGreen!65,fill=SPGreen!4] (info) at (0,-2.95)
 {\textbf{Información}\quad registros versionados · flooding · DAC/consenso};
\node[bar,draw=VIUOrange!75,fill=VIUOrange!4] (out) at (0,-4.15)
 {\textbf{Contrato físico}\quad asignación entera válida $a$ $\longrightarrow$ entrada de SP2};

\begin{pgfonlayer}{n4taxonomyback}
  \foreach \n in {f1,f2,f3,f4}{
    % Antes apuntaba a la esquina superior derecha, lo que producía una diagonal
    % suelta junto a F-IV. La flecha baja recta al centro de cada familia.
    \draw[arr] (problem.south -| \n.north) -- (\n.north);
    \draw[arr] (\n.south) -- (\n.south |- info.north);
  }
  \draw[arr] (info)--(out);
\end{pgfonlayer}
\end{tikzpicture}
}
\caption{Taxonomía de N4: variable, evolución, objetivo, salida física y nivel
de evidencia de las cuatro familias.}
\label{fig:n4-family-taxonomy}
\viuownsource
\end{figure}
\vspace{-3mm}

\clearpage

% --------------------------------------------------------------------------
% Atlas estratégico A - variantes atómicas
% --------------------------------------------------------------------------
\leveltag{RESULTADOS SP1 · N4 · ATLAS ESTRATÉGICO · F-I}
\section*{F-I: tres revisiones unilaterales sobre el mismo mundo}

Capacidades, demandas y agregados permanecen fijos en la
Figura~\ref{fig:n4-atlas-fi}; solo cambia la regla que selecciona la siguiente
acción. La mejor respuesta (BR) toma la mayor mejora; la revisión Smith atómica
asíncrona (ASR) sortea entre mejoras; y el aprendizaje log-lineal (LLL) conserva
una probabilidad positiva de aceptar un paso peor
\parencite{mardenShamma2012LogLinear}. La versión finita y bifásica evaluada
aquí no hereda automáticamente las garantías asintóticas de esa referencia. Las
tres reglas siguen operando en $\mathcal V_1$.

\begin{figure}[H]
\centering
\scalebox{.98}{\begin{tikzpicture}[x=1cm,y=1cm,>={Latex[length=1.25mm]},
 band/.style={draw=SPGrid,rounded corners=1.2mm,fill=SPFloor,
  minimum width=11.35cm,minimum height=2.30cm},
 robot/.style={draw=SPBlue,rounded corners=.55mm,fill=white,
  minimum width=.62cm,minimum height=.34cm,align=center,font=\fontsize{5.2}{5.5}\selectfont},
 load/.style={draw=VIUOrange,rounded corners=.55mm,fill=VIUOrange!5,
  minimum width=1.00cm,minimum height=.52cm,align=center,font=\fontsize{4.9}{5.2}\selectfont},
 rule/.style={draw=black!25,rounded corners=.9mm,fill=white,
  minimum width=2.10cm,minimum height=.82cm,align=center,font=\tiny\bfseries},
 assign/.style={->,line width=.6pt,draw=SPBlue,shorten <=1.1pt,shorten >=1.1pt},
 proposal/.style={->,densely dashed,line width=.8pt,draw=VIUOrange,shorten <=1.1pt,shorten >=1.1pt},
 accept/.style={->,line width=1.05pt,draw=SPGreen,shorten <=1.1pt,shorten >=1.1pt},
 explore/.style={->,densely dashed,line width=.85pt,draw=SPRed,shorten <=1.1pt,shorten >=1.1pt},
 flow/.style={->,thick,draw=black!40}]

\node[font=\scriptsize\bfseries,text=black!70] at (0,4.02)
 {$N=6$, $K=3$, $c_i=1$, $m=(2,2,2)$};

% --------------------------------------------------------------- BR
\node[band] (br) at (0,2.45) {};
\node[anchor=north west,font=\scriptsize\bfseries] at ($(br.north west)+(.17,-.13)$)
 {Geo-QPG-BR \textcolor{black!45}{$\mid$ $\arg\max_k g_i(k)$}};
\node[font=\tiny\bfseries,text=black!55] at (-3.82,2.92) {antes};
\node[robot] (br-r) at (-4.52,2.48) {$r_1$\\$c=1$};
\node[load] (br-k1) at (-4.62,1.78) {$k_1$\\$Q/m=3/2$};
\node[load] (br-k2) at (-3.52,1.78) {$k_2$\\$Q/m=1/2$};
\node[load] (br-k3) at (-2.42,1.78) {$k_3$\\$Q/m=2/2$};
\draw[assign] (br-r)--(br-k1);
\draw[proposal] (br-r) to[bend left=18] (br-k2);
\node[font=\tiny,fill=white,inner sep=.8pt] at (-3.34,2.30) {$g_i=0{,}8$};
\node[rule] (br-rule) at (0,2.30) {$k^\star=\arg\max_k g_i(k)$\\commit si $\Delta\Psi<_{\rm lex}0$};
\draw[flow] (-1.72,2.30)--(br-rule.west);
\draw[flow] (br-rule.east)--(1.72,2.30);
\node[font=\tiny\bfseries,text=black!55] at (3.82,2.92) {después};
\node[robot,draw=SPGreen] (br2-r) at (3.52,2.48) {$r_1$\\$c=1$};
\node[load] (br2-k1) at (2.42,1.78) {$k_1$\\$Q/m=2/2$};
\node[load] (br2-k2) at (3.52,1.78) {$k_2$\\$Q/m=2/2$};
\node[load] (br2-k3) at (4.62,1.78) {$k_3$\\$Q/m=2/2$};
\draw[accept] (br2-r)--(br2-k2);

% --------------------------------------------------------------- ASR
\node[band] (asr) at (0,.02) {};
\node[anchor=north west,font=\scriptsize\bfseries] at ($(asr.north west)+(.17,-.13)$)
 {Geo-ASR \textcolor{black!45}{$\mid$ $\Pr(k)\propto[g_i(k)]_+$}};
\node[font=\tiny\bfseries,text=black!55] at (-3.82,.49) {propuestas};
\node[robot] (as-r) at (-4.52,.05) {$r_1$\\$c=1$};
\node[load] (as-k1) at (-4.62,-.65) {$k_1$\\$Q/m=3/2$};
\node[load] (as-k2) at (-3.52,-.65) {$k_2$\\$Q/m=1/2$};
\node[load] (as-k3) at (-2.42,-.65) {$k_3$\\$Q/m=1/2$};
\draw[assign] (as-r)--(as-k1);
\draw[proposal,line width=1.15pt] (as-r) to[bend left=12] (as-k2);
\node[font=\tiny,fill=white,inner sep=.8pt] at (-3.52,-.30) {$P=0{,}6$};
\draw[proposal,line width=.55pt] (as-r) to[bend left=24] (as-k3);
\node[font=\tiny,fill=white,inner sep=.8pt] at (-2.42,-.30) {$P=0{,}4$};
\node[rule] (as-rule) at (0,-.13) {$P_i(k)=\dfrac{[g_i(k)]_+}{\sum_q[g_i(q)]_+}$\\misma vecindad $h=1$};
\draw[flow] (-1.72,-.13)--(as-rule.west);
\draw[flow] (as-rule.east)--(1.72,-.13);
\node[font=\tiny\bfseries,text=black!55] at (3.82,.49) {muestra};
\node[robot,draw=SPGreen] (as2-r) at (3.52,.05) {$r_1$\\$c=1$};
\node[load] (as2-k1) at (2.42,-.65) {$k_1$\\$Q/m=2/2$};
\node[load] (as2-k2) at (3.52,-.65) {$k_2$\\$Q/m=2/2$};
\node[load] (as2-k3) at (4.62,-.65) {$k_3$\\$Q/m=2/2$};
\draw[accept] (as2-r)--(as2-k2);

% --------------------------------------------------------------- LLL
\node[band] (lll) at (0,-2.41) {};
\node[anchor=north west,font=\scriptsize\bfseries] at ($(lll.north west)+(.17,-.13)$)
 {Geo-LLL \textcolor{black!45}{$\mid$ $\Pr(k)\propto e^{\beta u_i(k)}$}};
\node[font=\tiny\bfseries,text=black!55] at (-3.82,-1.94) {factible};
\node[robot] (ll-r) at (-3.52,-2.38) {$r_1$\\$c=1$};
\node[load] (ll-k1) at (-4.62,-3.08) {$k_1$\\$Q/m=2/2$};
\node[load] (ll-k2) at (-3.52,-3.08) {$k_2$\\$Q/m=2/2$};
\node[load] (ll-k3) at (-2.42,-3.08) {$k_3$\\$Q/m=2/2$};
\draw[assign] (ll-r)--(ll-k2);
\draw[explore] (ll-r) to[bend right=14] (ll-k1);
\node[font=\tiny,fill=white,inner sep=.8pt] at (-4.52,-2.56) {$P>0$};
\node[rule] (ll-rule) at (0,-2.56) {$P_i(k)\propto e^{\beta u_i(k)}$\\puede aceptar $\Delta\Psi>0$};
\draw[flow] (-1.72,-2.56)--(ll-rule.west);
\draw[flow] (ll-rule.east)--(1.72,-2.56);
\node[font=\tiny\bfseries,text=black!55] at (3.82,-1.94) {paso exploratorio};
\node[robot,draw=SPRed] (ll2-r) at (2.52,-2.38) {$r_1$\\$c=1$};
\node[load] (ll2-k1) at (2.42,-3.08) {$k_1$\\$Q/m=3/2$};
\node[load] (ll2-k2) at (3.52,-3.08) {$k_2$\\$Q/m=1/2$};
\node[load] (ll2-k3) at (4.62,-3.08) {$k_3$\\$Q/m=2/2$};
\draw[explore] (ll2-r)--(ll2-k1);

\end{tikzpicture}
}
\caption{Revisiones unilaterales de F-I sobre un mundo común. Los agregados
$Q_k/m_k$ incluyen los robots que permanecen fijos y aíslan el efecto del único
robot que cambia.}
\label{fig:n4-atlas-fi}
\viuownsource
\end{figure}

Cambiar la revisión amplía las acciones probadas, pero no el número de robots
que puede coordinarse. Por eso ASR y LLL no sustituyen una desviación de orden
dos o tres.

\clearpage

% --------------------------------------------------------------------------
% Atlas estratégico A2 - coordinación y scheduler
% --------------------------------------------------------------------------
\leveltag{RESULTADOS SP1 · N4 · ATLAS ESTRATÉGICO · COORDINACIÓN}
\section*{Orden local y ejecución concurrente}

El orden de una desviación y la ejecución concurrente son decisiones diferentes
(Figura~\ref{fig:n4-atlas-coalitional}). La mejor
respuesta bilateral (2BR) examina todas las desviaciones conjuntas de hasta dos
robots, incluidos los intercambios; C3 amplía la enumeración a grupos de hasta
tres, incluidos los ciclos. El conjunto independiente maximal distribuido con
\emph{commit} transaccional (DMIS+TX) recibe propuestas ya evaluadas, construye
el grafo de conflictos y confirma un conjunto independiente con versiones.

\begin{figure}[H]
\centering
\scalebox{.98}{\begin{tikzpicture}[x=1cm,y=1cm,>=Latex,
 band/.style={draw=SPGrid,rounded corners=1.2mm,fill=SPFloor,
  minimum width=11.35cm,minimum height=2.30cm},
 robot/.style={draw=SPBlue,rounded corners=.55mm,fill=white,
  minimum width=.62cm,minimum height=.34cm,align=center,font=\fontsize{5.2}{5.5}\selectfont},
 load/.style={draw=VIUOrange,rounded corners=.55mm,fill=VIUOrange!5,
  minimum width=.88cm,minimum height=.52cm,align=center,font=\fontsize{4.9}{5.2}\selectfont},
 rule/.style={draw=black!25,rounded corners=.9mm,fill=white,
  minimum width=2.10cm,minimum height=.82cm,align=center,font=\tiny\bfseries},
 proposal/.style={->,densely dashed,line width=.8pt,draw=VIUOrange},
 accept/.style={->,very thick,draw=SPGreen},
 flow/.style={->,thick,draw=black!40},
 prop/.style={circle,draw=SPPurple,fill=white,minimum size=4.8mm,font=\tiny}]

\node[font=\scriptsize\bfseries,text=black!70] at (0,4.02)
 {$N=6$, $K=3$, $c_i=1$, $m=(2,2,2)$};
\node[font=\tiny,text=black!55] at (0,3.70)
 {halo: AMR modificados \quad naranja: propuesta \quad verde: \emph{commit}};

% --------------------------------------------------------------- 2BR
\node[band] (pair) at (0,2.45) {};
\node[anchor=north west,font=\scriptsize\bfseries] at ($(pair.north west)+(.17,-.13)$)
 {Geo-QPG-2BR \textcolor{black!45}{$\mid$ $h=2$}};
\node[font=\tiny\bfseries,text=black!55] at (-3.82,2.92) {bloqueo $1$-local};
\node[robot] (p-r1) at (-4.35,2.48) {$r_1$};
\node[robot] (p-r2) at (-3.20,2.48) {$r_2$};
\node[load] (p-k1) at (-4.35,1.78) {$k_1$\\$Q/m=2/2$};
\node[load] (p-k2) at (-3.20,1.78) {$k_2$\\$Q/m=2/2$};
\node[load] (p-k3) at (-2.10,1.78) {$k_3$\\$Q/m=2/2$};
\draw[SPPurple,dashed,thick] (-4.02,2.42) ellipse (1.12 and .42);
\draw[proposal] (p-r1) to[bend left=18] (p-k2);
\draw[proposal] (p-r2) to[bend right=18] (p-k1);
\node[rule] (p-rule) at (0,2.30) {$\Delta D=0$\\$\Delta J<0$};
\draw[flow] (-1.55,2.30)--(p-rule.west); \draw[flow] (p-rule.east)--(1.55,2.30);
\node[font=\tiny\bfseries,text=black!55] at (3.82,2.92) {perfil aceptado};
\node[robot,draw=SPGreen] (p2-r1) at (3.20,2.48) {$r_1$};
\node[robot,draw=SPGreen] (p2-r2) at (4.35,2.48) {$r_2$};
\node[load] (p2-k1) at (3.20,1.78) {$k_1$\\$Q/m=2/2$};
\node[load] (p2-k2) at (4.35,1.78) {$k_2$\\$Q/m=2/2$};
\node[load] (p2-k3) at (2.10,1.78) {$k_3$\\$Q/m=2/2$};
\draw[accept] (p2-r1)--(p2-k1); \draw[accept] (p2-r2)--(p2-k2);

% --------------------------------------------------------------- C3
\node[band] (triple) at (0,.02) {};
\node[anchor=north west,font=\scriptsize\bfseries] at ($(triple.north west)+(.17,-.13)$)
 {Geo-QPG-C3 \textcolor{black!45}{$\mid$ $h=3$}};
\node[font=\tiny\bfseries,text=black!55] at (-3.82,.49) {bloqueo $2$-local};
\node[robot] (c-r1) at (-4.52,.05) {$r_1$};
\node[robot] (c-r2) at (-3.52,.05) {$r_2$};
\node[robot] (c-r3) at (-2.52,.05) {$r_3$};
\node[load] (c-k1) at (-4.52,-.65) {$k_1$\\$Q/m=2/2$};
\node[load] (c-k2) at (-3.52,-.65) {$k_2$\\$Q/m=2/2$};
\node[load] (c-k3) at (-2.52,-.65) {$k_3$\\$Q/m=2/2$};
\draw[SPPurple,dashed,thick] (-3.52,-.02) ellipse (2.05 and .32);
\draw[proposal] (c-r1) to[bend left=15] (c-k2);
\draw[proposal] (c-r2) to[bend left=15] (c-k3);
\draw[proposal] (c-r3) to[bend right=10] (c-k1);
\node[rule] (c-rule) at (0,-.13) {$r_1:1\to2$\\$r_2:2\to3$\\$r_3:3\to1$};
\draw[flow] (-1.55,-.13)--(c-rule.west); \draw[flow] (c-rule.east)--(1.55,-.13);
\node[font=\tiny\bfseries,text=black!55] at (3.82,.49) {desviación aceptada};
\node[robot,draw=SPGreen] (c2-r1) at (3.52,.05) {$r_1$};
\node[robot,draw=SPGreen] (c2-r2) at (4.52,.05) {$r_2$};
\node[robot,draw=SPGreen] (c2-r3) at (2.52,.05) {$r_3$};
\node[load] (c2-k1) at (2.52,-.65) {$k_1$\\$Q/m=2/2$};
\node[load] (c2-k2) at (3.52,-.65) {$k_2$\\$Q/m=2/2$};
\node[load] (c2-k3) at (4.52,-.65) {$k_3$\\$Q/m=2/2$};
\draw[accept] (c2-r1)--(c2-k2); \draw[accept] (c2-r2)--(c2-k3); \draw[accept] (c2-r3)--(c2-k1);

% --------------------------------------------------------------- DMIS
\node[band] (dmis) at (0,-2.41) {};
\node[anchor=north west,font=\scriptsize\bfseries] at ($(dmis.north west)+(.17,-.13)$)
 {Geo-QPG-DMIS+TX \textcolor{black!45}{$\mid$ concurrencia segura}};
\node[font=\tiny\bfseries,text=black!55] at (-3.82,-1.94) {propuestas vecinales};
\node[prop] (d-p1) at (-4.52,-2.48) {$p_1$};
\node[prop] (d-p2) at (-3.52,-2.48) {$p_2$};
\node[prop] (d-p3) at (-2.52,-2.48) {$p_3$};
\node[font=\tiny,text=black!55] at (-3.52,-3.05) {cada $p$ declara AMR, cargas y versiones};
\node[rule] (d-rule) at (0,-2.56) {grafo de conflictos\\DMIS: $\{p_1,p_3\}$\\PREPARE / COMMIT};
\draw[flow] (-1.55,-2.56)--(d-rule.west); \draw[flow] (d-rule.east)--(1.55,-2.56);
\draw[SPPurple,thick] (d-p1)--(d-p2)--(d-p3);
\node[prop,draw=SPGreen,fill=SPGreen!8] (d2-p1) at (2.52,-2.48) {$p_1$};
\node[prop,draw=black!25] (d2-p2) at (3.52,-2.48) {$p_2$};
\node[prop,draw=SPGreen,fill=SPGreen!8] (d2-p3) at (4.52,-2.48) {$p_3$};
\node[font=\tiny\bfseries,text=black!55] at (3.52,-1.94) {propuestas compatibles};
\draw[accept] (d2-p1)--++(0,-.43); \draw[accept] (d2-p3)--++(0,-.43);
\node[font=\tiny,text=black!55] at (3.52,-3.05) {$v_k:4\to5$; $p_2$ se reevalúa};

\end{tikzpicture}
}
\caption{Extensiones coalicionales y transaccionales de F-I sobre el mismo
mundo. El halo identifica los robots que deben cambiar juntos; el último panel
opera sobre propuestas, no sobre una utilidad distinta.}
\label{fig:n4-atlas-coalitional}
\viuownsource
\end{figure}

La distinción será visible en resultados: aumentar $h$ puede mejorar el mínimo
local; DMIS+TX cambia concurrencia y tráfico. Ninguna de las dos propiedades se
deduce de la otra.

\clearpage

% --------------------------------------------------------------------------
% Atlas estratégico B - formulaciones continuas y temporal
% --------------------------------------------------------------------------
\leveltag{RESULTADOS SP1 · N4 · ATLAS ESTRATÉGICO · F-II--F-IV}
\section*{Variables estratégicas y salida física}

Tres variables estratégicas distintas pueden referirse al mismo mundo
(Figura~\ref{fig:n4-atlas-continuous}). F-II mueve intención sobre un símplex; F-III
añade precios de las cuotas compartidas; F-IV reacciona a una secuencia de
eventos. Solo $a$ puede entregarse directamente a SP2. Los estados continuos
pasan por el mismo cierre $\mathcal R$, y la rama temporal produce una secuencia
de asignaciones $a_t$.

\begin{figure}[H]
\centering
\scalebox{.98}{\begin{tikzpicture}[x=1cm,y=1cm,>=Latex,
 band/.style={draw=SPGrid,rounded corners=1.2mm,fill=SPFloor,
  minimum width=11.40cm,minimum height=2.48cm},
 card/.style={draw=black!23,rounded corners=.9mm,fill=white,
  minimum width=2.35cm,minimum height=1.18cm,align=center,font=\tiny},
 stage/.style={draw=black!23,rounded corners=.9mm,fill=white,
  minimum width=2.12cm,minimum height=1.18cm,align=center,font=\tiny},
 arrow/.style={->,thick,draw=VIUOrange},
 down/.style={->,thick,draw=black!40}]

\node[font=\scriptsize\bfseries,text=black!70] at (0,4.80)
 {MISMO MUNDO FÍSICO; TRES ESTADOS ESTRATÉGICOS};
\node[font=\tiny,text=black!55] at (0,4.49)
 {$a$ es ejecutable; $x$ y $(x,\lambda)$ pasan por $\mathcal R$; F-IV produce una secuencia $a_t$};

% ================================================================ F-II
\node[band] (fii) at (0,3.12) {};
\node[anchor=north west,font=\scriptsize\bfseries,text=SPBlue]
 at ($(fii.north west)+(.18,-.13)$) {F-II · INTENCIÓN CONTINUA $x_i\in\Delta_i$};
\node[card] (rep) at (-4.12,3.14)
 {\textbf{Replicator}\\[.6mm]
  $x_i=(.7,.3,0)$\\
  una componente en cero\\
  no se reactiva};
\node[card] (smith) at (-1.38,3.14)
 {\textbf{Smith}\\[.6mm]
  $k_1\rightleftarrows k_2$\\
  transfiere masa según\\
  diferencias de fitness};
\node[card] (bnn) at (1.38,3.14)
 {\textbf{BNN}\\[.6mm]
  $\widehat F_{ik}=[F_{ik}-\bar F_i]_+$\\
  entra masa donde el\\
  exceso es positivo};
\node[card] (logit) at (4.12,3.14)
 {\textbf{Logit}\\[.6mm]
  $x_i^+\propto\exp(F_i/\tau)$\\
  softmax difuso; $\tau$\\
  regula exploración};
\node[draw=VIUOrange,rounded corners=.8mm,fill=VIUOrange!4,
 minimum width=4.15cm,minimum height=.46cm,font=\tiny\bfseries] (closure) at (0,2.05)
 {$x^\star\quad\xrightarrow{\ \mathcal R\ }\quad a^\star$ entero};
\foreach \n in {rep,smith,bnn,logit}{\draw[down] (\n.south)--(closure.north);}
\node[font=\fontsize{5.3}{5.6}\selectfont,text=SPRed] at (0,1.53)
 {las cuatro dinámicas comparten $F$ y el cierre; el ranking puede cambiar después de $\mathcal R$};

% ================================================================ F-III
\node[band] (fiii) at (0,.00) {};
\node[anchor=north west,font=\scriptsize\bfseries,text=SPGreen!70!black]
 at ($(fiii.north west)+(.18,-.13)$) {F-III · CUOTAS COMO RESTRICCIÓN COMPARTIDA};
\node[stage,draw=VIUOrange!70] (scarcity) at (-4.18,.02)
 {\textbf{1 · escasez}\\[.6mm]
  carga $k_1$\\
  $Q_1/m_1=2/5$\\
  $g_1(x)>0$};
\node[stage,draw=SPRed!60] (price) at (-1.40,.02)
 {\textbf{2 · precio}\\[.6mm]
  $\lambda_1\uparrow$\\
  $c_i\lambda_1-d_{i1}$\\
  atrae capacidad};
\node[stage,draw=SPPurple!60] (consensus) at (1.40,.02)
 {\textbf{3 · consenso}\\[.6mm]
  $\lambda_{11}\leftrightarrow\lambda_{21}$\\
  $\leftrightarrow\lambda_{31}$\\
  copias vecinales};
\node[stage,draw=SPGreen!65] (physical) at (4.18,.02)
 {\textbf{4 · salida}\\[.6mm]
  $(x^\star,\lambda^\star)$\\
  $\mathcal R(x^\star)=a^\star$\\
  coalición atómica};
\draw[arrow] (scarcity.east)--(price.west);
\draw[arrow] (price.east)--(consensus.west);
\draw[arrow] (consensus.east)--(physical.west);
\node[font=\tiny,text=black!55] at (0,-1.00)
 {aceptación numérica: residual primal + dual + complementariedad + estación + consenso};

% ================================================================ F-IV
\node[band] (fiv) at (0,-3.12) {};
\node[anchor=north west,font=\scriptsize\bfseries,text=SPPurple]
 at ($(fiv.north west)+(.18,-.13)$) {F-IV · RECLUTAMIENTO BAJO EVENTOS};
\node[stage] (t0) at (-4.18,-3.10)
 {\textbf{$t_0$ · servicio}\\[.6mm]
  cargas activas\\
  robots disponibles\\
  coalición $a_t$};
\node[stage,draw=VIUOrange!65] (t1) at (-1.40,-3.10)
 {\textbf{$t_1$ · llegada}\\[.6mm]
  aparece $k_{\rm new}$\\
  cambia la demanda\\
  estado $s_{t+1}$};
\node[stage,draw=SPRed!65] (t2) at (1.40,-3.10)
 {\textbf{$t_2$ · fallo}\\[.6mm]
  $r_j$ deja de estar\\
  disponible\\
  se invalida parte de $a_t$};
\node[stage,draw=SPGreen!65] (t3) at (4.18,-3.10)
 {\textbf{$t_3$ · política}\\[.6mm]
  reactiva / inicialización previa\\
  oráculo escalar finito\\
  nueva asignación $a_{t+1}$};
\draw[arrow] (t0.east)--node[above,font=\tiny] {evento}(t1.west);
\draw[arrow] (t1.east)--node[above,font=\tiny] {$T(s,u,e)$}(t2.west);
\draw[arrow] (t2.east)--node[above,font=\tiny] {decisión}(t3.west);
\node[font=\tiny,text=SPRed] at (0,-4.14)
 {estudio temporal con eventos exógenos; sin garantía de potencial de Markov};

\end{tikzpicture}
}
\caption{Atlas estratégico de F-II--F-IV: estado, actualización y conversión a una
salida atómica. La comparación estática termina en $a^\star$; la temporal se
evalúa sobre la secuencia completa.}
\label{fig:n4-atlas-continuous}
\viuownsource
\end{figure}

El atlas fija también qué significa ``converger''. En F-II es un criterio
numérico de la dinámica y del estimador; en F-III deben cruzarse conjuntamente
los residuos de Karush--Kuhn--Tucker (KKT) y de consenso. En F-IV no hay un
equilibrio estático único:
se compara coste acumulado, servicio y cambios de coalición bajo los mismos
eventos.

\clearpage

% --------------------------------------------------------------------------
% Página 28 - F-I, potencial y suficiencia marginal
% --------------------------------------------------------------------------
\leveltag{RESULTADOS SP1 · N4 · F-I · JUEGO POTENCIAL ATÓMICO}
\section*{F-I: una mejora local no necesita el perfil completo}
\enlargethispage{2\baselineskip}

Geo-QPG opera sobre $a$ y evita sumar kg y metros con un peso arbitrario. Compara
\begin{equation}
\label{eq:n4-state-v2}
\Psi_4(a)=\bigl(D_4(a),J_4(a)\bigr)
\end{equation}
en orden lexicográfico: primero reduce déficit y, con $D_4=0$, conserva cobertura
y reduce distancia. El MILP de N2 obtiene el mínimo global; F-I busca descensos
locales con los registros que toca cada propuesta.

Sea $p:i:r\rightarrow s$ una desviación unilateral y $d_{i0}=0$. Al omitir la
carga inactiva, la variación exacta es
\begin{equation}
\label{eq:n4-unilateral-delta}
\begin{aligned}
\Delta_iD={}&[m_r-(Q_r-c_i^{\mathrm{pay}})]_+
 +[m_s-(Q_s+c_i^{\mathrm{pay}})]_+\\[-1mm]
&-[m_r-Q_r]_+-[m_s-Q_s]_+,\\
\Delta_iJ={}&d_{is}-d_{ir}.
\end{aligned}
\end{equation}

Para valorar $i:r\rightarrow s$ bastan $Q_r,Q_s,c_i^{\mathrm{pay}},d_{ir}$ y
$d_{is}$. La identidad se derivó algebraicamente y se comprobó por enumeración.
Una propuesta coalicional añade únicamente sus participantes y las cargas que
modifica.

\begin{figure}[H]
\centering
\scalebox{.88}{\begin{tikzpicture}[x=1cm,y=1cm,>=Latex,
 load/.style={draw=VIUOrange!85,rounded corners=1mm,fill=VIUOrange!5,
  minimum width=2.35cm,minimum height=1.35cm,align=center,font=\footnotesize},
 reg/.style={draw=SPGreen!70,rounded corners=.8mm,fill=SPGreen!5,
  minimum width=1.65cm,minimum height=.62cm,align=center,font=\scriptsize},
 own/.style={draw=SPBlue!70,rounded corners=.8mm,fill=SPBlue!5,
  minimum width=2.55cm,minimum height=.82cm,align=center,font=\scriptsize},
 robot/.style={circle,draw=SPBlue,fill=white,minimum size=6.4mm,font=\scriptsize},
 arr/.style={->,very thick,draw=VIUOrange}]
\node[load] (r) at (-3.75,.65) {carga $r$\\demanda $m_r$};
\node[load] (s) at (3.75,.65) {carga $s$\\demanda $m_s$};
\node[reg] (qr) at (-3.75,-.52) {$z_r=(Q_r,v_r)$};
\node[reg] (qs) at (3.75,-.52) {$z_s=(Q_s,v_s)$};
\node[robot] (i) at (-1.82,.65) {$i$};
\node[own] (state) at (0,-1.05)
 {estado propio: $c_i^{\mathrm{pay}},\ d_{ir},\ d_{is}$};
\node[draw=SPPurple!70,rounded corners=1mm,fill=SPPurple!5,
 minimum width=2.5cm,minimum height=.9cm,align=center,font=\footnotesize] (delta) at (0,1.55)
 {evaluación exacta\\$(\Delta_iD,\Delta_iJ)$};
\draw[arr] (i) to[bend right=10] node[below,font=\scriptsize] {$i:r\rightarrow s$} (s);
\draw[->,thin,draw=SPGreen!80] (qr)--(delta);
\draw[->,thin,draw=SPGreen!80] (qs)--(delta);
\draw[->,thin,draw=SPBlue!80] (state)--(delta);
\node[font=\scriptsize,align=center,text=black!55] at (0,-1.85)
 {las cargas no afectadas se cancelan; las versiones sí deben llegar por la red};
\end{tikzpicture}
}
\caption{Una desviación se evalúa con estado propio y registros versionados de
las cargas afectadas; el resto del perfil se cancela.}
\label{fig:n4-sufficient-information}
\viuownsource
\end{figure}

Los agregados y sus versiones aún deben llegar por la red. Con $\Phi_Q=-D_4$ y,
sobre $\mathcal F=\{a:D_4(a)=0\}$, $\Phi_G=-J_4$, se define
\begin{equation}
\label{eq:n4-wonderful-life}
u_i^X(a)=\Phi_X(a)-\Phi_X(0,a_{-i}),\qquad X\in\{Q,G\}.
\end{equation}

\begin{n1prop}[Potencial exacto unilateral]
Dentro de una misma fase,
$u_i^X(a_i',a_{-i})-u_i^X(a_i,a_{-i})
=\Phi_X(a_i',a_{-i})-\Phi_X(a_i,a_{-i})$.
\end{n1prop}
El término de referencia se cancela~\parencite{monderer1996potential}. Para
$h>1$ se hablará de búsqueda finita coalicional, no de potencial exacto.

\clearpage

% --------------------------------------------------------------------------
% Página 29 - jerarquía h-local y reglas de revisión
% --------------------------------------------------------------------------
\leveltag{RESULTADOS SP1 · N4 · F-I · ORDEN ESTRATÉGICO Y REVISIÓN}
\section*{Orden de coordinación y regla de revisión}
\enlargethispage{4\baselineskip}

La vecindad
\begin{equation}
\label{eq:n4-h-neighbourhood}
\mathcal V_h(a)=\left\{a':1\leq
\left|\{i:a_i'\neq a_i\}\right|\leq h\right\}
\end{equation}
contiene los perfiles alcanzables coordinando hasta $h$ robots. Un perfil es
\emph{$h$-local estable} si ningún $a'\in\mathcal V_h(a)$ reduce
$\Psi_4$. Así, BR, 2BR y C3 son las etapas $h=1,2,3$ de una sola jerarquía,
con $\mathcal V_1\subset\mathcal V_2\subset\mathcal V_3$.

La regla de revisión responde a otra pregunta: cuál de las opciones disponibles
se prueba. Para la ganancia unilateral $g_i^X(k)$, las variantes ejecutadas
fueron
\begin{equation}
\label{eq:n4-revision-rules}
\begin{aligned}
\mathrm{BR}:\;&k^\star\in\arg\max_k g_i^X(k),\quad
g_i^X(k^\star)>0,\\[-1mm]
\mathrm{Geo\mbox{-}ASR}:\;&
\Pr(k)=\frac{[g_i^X(k)]_+}{\sum_q[g_i^X(q)]_+},\\[-1mm]
\mathrm{Geo\mbox{-}LLL}:\;&
\Pr(k)=\frac{\exp\{\beta u_i^X(k,a_{-i})\}}
{\sum_q\exp\{\beta u_i^X(q,a_{-i})\}}.
\end{aligned}
\end{equation}
BR no publica una propuesta si la mejor ganancia no es positiva. En ASR, si
$\sum_q[g_i^X(q)]_+=0$, el robot conserva $a_i$ y no se evalúa la razón.

\begin{figure}[H]
\centering
\begin{tikzpicture}[x=1cm,y=1cm,>=Latex,
 panel/.style={draw=SPGrid,rounded corners=1mm,fill=SPFloor,
  text width=3.05cm,minimum height=1.70cm,align=center,font=\scriptsize},
 rob/.style={circle,draw=SPBlue,fill=white,minimum size=5.2mm,font=\tiny},
 load/.style={draw=VIUOrange,rounded corners=.6mm,fill=VIUOrange!5,
  minimum width=6mm,minimum height=5mm,font=\tiny},
 rule/.style={draw=SPGrid,rounded corners=1mm,text width=3.05cm,
  minimum height=1.52cm,align=center,inner sep=1.6mm,font=\scriptsize}]
\node[font=\footnotesize\bfseries,anchor=west] at (-5.25,2.15)
 {A · orden estratégico: cuántos robots pueden cambiar juntos};
\foreach \x/\title/\h in {-3.75/{BR}/1,0/{2BR}/2,3.75/{C3}/3}{
  \node[panel] (p\h) at (\x,.95) {};
  \node[font=\scriptsize\bfseries] at ($(p\h.center)+(0,.58)$) {\title\quad $h=\h$};
  \node[load] at ($(p\h.center)+(-.9,-.48)$) {$k_1$};
  \node[load] at ($(p\h.center)+(.9,-.48)$) {$k_2$};
  \foreach \j in {1,...,3}{\node[rob] at ($(p\h.center)+(-.7+.7*\j,.08)$) {$\j$};}
}
\draw[very thick,draw=SPBlue] ($(p1.center)+(-.02,.08)$) circle (3.1mm);
\draw[very thick,draw=SPGreen] ($(p2.center)+(.35,.08)$) ellipse (7.4mm and 3.6mm);
\draw[very thick,draw=VIUOrange] ($(p3.center)+(.7,.08)$) ellipse (10.8mm and 3.6mm);
\draw[->,thick,draw=VIUOrange] ($(p1.east)+(0,.48)$)--node[above,font=\tiny] {$\mathcal V_1(a)\subseteq\mathcal V_2(a)$}($(p2.west)+(0,.48)$);
\draw[->,thick,draw=VIUOrange] ($(p2.east)+(0,.48)$)--node[above,font=\tiny] {$\mathcal V_2(a)\subseteq\mathcal V_3(a)$}($(p3.west)+(0,.48)$);

\node[font=\footnotesize\bfseries,anchor=west] at (-5.25,-.40)
 {B · regla de revisión: cómo se elige dentro de la misma $\mathcal V_1(a)$};
\node[rule,draw=SPBlue!70,fill=SPBlue!4] (br) at (-3.75,-1.65)
 {\textbf{BR}\\elige $\arg\max_k g_i(k)$\\si $g_i(k)>0$};
\node[rule,draw=SPGreen!70,fill=SPGreen!4] (asr) at (0,-1.65)
 {\textbf{Geo-ASR}\\sortea mejoras\\con peso $[g_i(k)]_+$};
\node[rule,draw=VIUOrange!80,fill=VIUOrange!4] (lll) at (3.75,-1.65)
 {\textbf{Geo-LLL}\\respuesta de Gibbs;\\admite peores\\transitorias};
\node[font=\scriptsize,text=black!55] at (0,-2.72)
 {cambiar la regla explora de otra forma; no aumenta el número de robots coordinados};
\end{tikzpicture}

\caption{El orden estratégico amplía la vecindad; BR, Geo-ASR y Geo-LLL
cambian cómo se explora la misma vecindad unilateral.}
\label{fig:n4-order-revision}
\viuownsource
\end{figure}

Geo-ASR es una revisión atómica de tipo Smith, no la dinámica continua de F-II.
Geo-LLL explora durante $\min\{8N,\lfloor T/2\rfloor\}$ activaciones con
$\beta_t:1\to24$ y después aplica BR. En cobertura admite pasos exploratorios;
tras alcanzar $D_4=0$, solo admite acciones que lo conservan. El terminal es
1-local solo si el pulido BR no queda truncado. Ninguna
de las dos coordina por sí misma a dos robots; por tanto, no puede convertir un
bloqueo bilateral en una mejora unilateral. Esta distinción fija la lectura del
experimento E4: primero se compara la revisión dentro de $h=1$ y luego el orden
$h$.

\clearpage

% --------------------------------------------------------------------------
% Página 30 - barrera y refinamiento progresivo
% --------------------------------------------------------------------------
\leveltag{RESULTADOS SP1 · N4 · F-I · BARRERAS Y REFINAMIENTO}
\section*{Orden de escape y refinamiento de la vecindad}
\enlargethispage{1\baselineskip}
\begingroup
\setlength{\abovedisplayskip}{5pt}
\setlength{\belowdisplayskip}{5pt}
\setlength{\abovedisplayshortskip}{3pt}
\setlength{\belowdisplayshortskip}{3pt}

La definición general no impone conectividad:
\begin{equation}
\label{eq:n4-escape-order}
h^\star(a)=\min\{h:\exists a'\in\mathcal V_h(a),\
\Psi_4(a')<_{\mathrm{lex}}\Psi_4(a)\},
\end{equation}
con $h^\star=\infty$ cuando no existe mejora. E9 audita una cantidad más
restrictiva. Para el grafo vecinal $G$, sea
\begin{equation}
\label{eq:n4-connected-escape-order}
\begin{aligned}
\mathcal V_h^{\mathrm c}(a)&=\{a'\in\mathcal V_h(a):
G[\{i:a_i'\ne a_i\}]\ \text{es conexo}\},\\[-1mm]
h_{\mathrm c}^\star(a)&=\min\{h:\exists a'\in\mathcal V_h^{\mathrm c}(a),\
\Psi_4(a')<_{\mathrm{lex}}\Psi_4(a)\}.
\end{aligned}
\end{equation}
Para dos cargas $r,s$ y dos
robots $i,j$, una barrera unilateral con escape bilateral aparece si
\begin{equation}
\label{eq:n4-bilateral-barrier-compact}
\begin{gathered}
Q_r-c_i^{\mathrm{pay}}<m_r,\quad Q_s-c_j^{\mathrm{pay}}<m_s,\\
Q_r-c_i^{\mathrm{pay}}+c_j^{\mathrm{pay}}\ge m_r,\quad
Q_s-c_j^{\mathrm{pay}}+c_i^{\mathrm{pay}}\ge m_s,\\
d_{is}+d_{jr}<d_{ir}+d_{js}.
\end{gathered}
\end{equation}

\begin{n1lema}[Barrera unilateral y escape bilateral]
Bajo~\eqref{eq:n4-bilateral-barrier-compact}, retirar solo a $i$ o a $j$ crea
déficit, mientras el intercambio mantiene $D_4=0$ y hace $\Delta J_4<0$. En el
subproblema de esas dos cargas, $h^\star(a)=2$.
\end{n1lema}

\begin{figure}[H]
\centering
\scalebox{.74}{\begin{tikzpicture}[x=1cm,y=1cm,>=Latex,
 panel/.style={draw=SPGrid,rounded corners=1mm,fill=SPFloor,
  minimum width=2.55cm,minimum height=2.12cm,align=center,font=\scriptsize},
 load/.style={draw=VIUOrange,rounded corners=.6mm,fill=VIUOrange!6,
  minimum width=7mm,minimum height=5.5mm,font=\tiny},
 rob/.style={circle,draw=SPBlue,fill=white,minimum size=5.3mm,font=\tiny},
 good/.style={->,very thick,draw=SPGreen},bad/.style={->,thick,draw=SPRed,dashed}]
\foreach \x/\name/\title in {-4.35/a/{A · bloqueado},-1.45/b/{B · intento unilateral},1.45/c/{C · propuesta bilateral},4.35/d/{D · perfil aceptado}}{
  \node[panel] (\name) at (\x,.75) {};
  % A \scriptsize los títulos medían más que el panel (2,55 cm) e invadían el vecino.
  \node[font=\tiny\bfseries,anchor=north]
    at ($(\name.center)+(0,.93)$) {\title};
  \node[load] (\name-r) at ($(\name.center)+(-.65,-.35)$) {$r$};
  \node[load] (\name-s) at ($(\name.center)+(.65,-.35)$) {$s$};
}
\node[rob] (ai) at ($(a.center)+(-.65,.35)$) {$i$};
\node[rob] (aj) at ($(a.center)+(.65,.35)$) {$j$};
\draw[->,thin,draw=SPBlue] (ai)--(a-r); \draw[->,thin,draw=SPBlue] (aj)--(a-s);
\node[rob] (bi) at ($(b.center)+(-.65,.35)$) {$i$};
\node[rob] (bj) at ($(b.center)+(.65,.35)$) {$j$};
\draw[bad] (bi) to[bend left=18] (b-s);
\node[font=\tiny,fill=white,inner sep=.7pt] at (-1.46,.76) {$D>0$};
\draw[->,thin,draw=SPBlue] (bj)--(b-s);
\node[rob] (ci) at ($(c.center)+(-.65,.35)$) {$i$};
\node[rob] (cj) at ($(c.center)+(.65,.35)$) {$j$};
\draw[good] (ci) to[bend left=18] (c-s);
\draw[good] (cj) to[bend right=18] (c-r);
\node[font=\tiny,text=SPGreen!80!black] at ($(c.center)+(0,-.83)$) {$\Delta D=0,\ \Delta J<0$};
\node[rob] (di) at ($(d.center)+(.65,.35)$) {$i$};
\node[rob] (dj) at ($(d.center)+(-.65,.35)$) {$j$};
\draw[->,thin,draw=SPGreen] (di)--(d-s); \draw[->,thin,draw=SPGreen] (dj)--(d-r);
\draw[->,thick,draw=black!50] (a.east)--(b.west);
\draw[->,thick,draw=black!50] (b.east)--(c.west);
\draw[->,thick,draw=black!50] (c.east)--(d.west);

\node[draw=SPBlue!65,rounded corners=.8mm,fill=SPBlue!4,
 minimum width=2.55cm,minimum height=.62cm,font=\scriptsize] (h1) at (-3.4,-1.25)
 {$a^{(1)}$: agotar $h=1$};
\node[draw=SPGreen!65,rounded corners=.8mm,fill=SPGreen!4,
 minimum width=2.55cm,minimum height=.62cm,font=\scriptsize] (h2) at (0,-1.25)
 {$a^{(2)}$: refinar con $h=2$};
\node[draw=VIUOrange!75,rounded corners=.8mm,fill=VIUOrange!4,
 minimum width=2.55cm,minimum height=.62cm,font=\scriptsize] (h3) at (3.4,-1.25)
 {$a^{(3)}$: refinar con $h=3$};
\draw[->,thick,draw=VIUOrange] (h1)--(h2);
\draw[->,thick,draw=VIUOrange] (h2)--(h3);
\node[font=\tiny,text=black!55] at (0,-1.98)
 {$\Psi_4(a^{(3)})\leq_{\rm lex}\Psi_4(a^{(2)})\leq_{\rm lex}\Psi_4(a^{(1)})$ bajo mejora estricta y búsqueda justa};
\end{tikzpicture}
}
\caption{Barrera bilateral y refinamiento progresivo entre etapas.}
\label{fig:n4-bilateral-barrier}
\viuownsource
\end{figure}
\vspace{-4mm}

\begin{algorithm}[H]
\caption{Refinamiento $h$-local de Geo-QPG}
\label{alg:n4-progressive}
\footnotesize
\begin{algorithmic}[1]
\Require perfil $a$ y orden máximo $h_{\max}$
\For{$h=1,\ldots,h_{\max}$}
  \State aceptar mejoras estrictas de $\mathcal V_h(a)$ hasta agotarlas o alcanzar el límite temporal
  \State $a^{(h)}\gets a$
\EndFor
\State \Return $a^{(h_{\max})}$
\end{algorithmic}
\end{algorithm}
\vspace{-2mm}

\begin{n1prop}[Terminación sin truncamiento]
Si se examinan todos los candidatos, cada \emph{commit} reduce $\Psi_4$ y no hay
interrupción por presupuesto, el espacio finito de $(K+1)^N$ perfiles garantiza
la terminación en un mínimo $h$-local.
\end{n1prop}
Una ejecución truncada no certifica estabilidad $h$-local ni la terminación
implica optimalidad global. E9 mide $h_{\mathrm c}^\star$ en el terminal
unilateral, con búsqueda exacta solo hasta $h=3$.
\endgroup

\clearpage

% --------------------------------------------------------------------------
% F-I - conflicto, DMIS y commit
% --------------------------------------------------------------------------
\leveltag{RESULTADOS SP1 · N4 · F-I · CONCURRENCIA Y TRANSACCIÓN}
\section*{De la mejora local a una ejecución concurrente}
\enlargethispage{1\baselineskip}

Sea $\mathcal P$ el conjunto de propuestas de una misma fase y versión. El grafo
$\mathcal G_C=(\mathcal P,\mathcal E_C)$ enlaza las que comparten robot o carga.

\begin{n1prop}[Composición de propuestas no conflictivas]
Si $\mathcal B\subseteq\mathcal P$ es independiente, entonces
$\Delta D_{\mathcal B}=\sum_{p\in\mathcal B}\Delta_pD$ y
$\Delta J_{\mathcal B}=\sum_{p\in\mathcal B}\Delta_pJ$; si cada propuesta mejora,
el lote también reduce $\Psi_4$.
\end{n1prop}
La independencia garantiza composición segura, aunque no selecciona
necesariamente el mejor lote.

\begin{figure}[H]
\centering
\scalebox{.74}{\begin{tikzpicture}[x=1cm,y=1cm,>=Latex,
 load/.style={draw=VIUOrange,rounded corners=.6mm,fill=VIUOrange!5,
  minimum width=8mm,minimum height=6mm,font=\scriptsize},
 prop/.style={circle,draw=SPBlue,fill=white,minimum size=7mm,font=\scriptsize},
 chosen/.style={circle,draw=SPGreen!80!black,fill=SPGreen!10,very thick,
  minimum size=7mm,font=\scriptsize},
 arr/.style={->,thick,draw=VIUOrange}]
\node[font=\footnotesize\bfseries] at (-3.4,1.75) {propuestas de reasignación};
\node[load] (k1) at (-4.45,.55) {$k_1$};
\node[load] (k2) at (-2.35,.55) {$k_2$};
\node[prop] (rp1) at (-4.45,-.7) {$p_1$};
\node[prop] (rp2) at (-3.55,-.7) {$p_2$};
\node[prop] (rp3) at (-2.35,-.7) {$p_3$};
\node[prop] (rp4) at (-1.45,-.7) {$p_4$};
\draw[arr] (rp1)--(k1); \draw[arr] (rp2)--(k1);
\draw[arr] (rp3)--(k2); \draw[arr] (rp4)--(k2);
% La etiqueta iba justo bajo el segmento y quedaba sobre los círculos p2 y p3.
\draw[dashed,draw=SPRed,thick] (rp2)--(rp3);
\node[font=\tiny,fill=white,inner sep=.7pt,below=3.2mm] at ($(rp2)!.5!(rp3)$)
  {AMR compartido};

\node[font=\footnotesize\bfseries] at (3.2,1.75) {grafo de conflictos $\mathcal G_C$};
\node[chosen] (p1) at (1.35,.65) {$p_1$};
\node[prop] (p2) at (3.0,1.1) {$p_2$};
\node[prop] (p3) at (3.0,-.25) {$p_3$};
\node[chosen] (p4) at (4.65,.65) {$p_4$};
\draw[thick,draw=SPRed] (p1)--(p2)--(p3)--(p4);
\draw[thick,draw=SPRed] (p2)--(p4);
\node[draw=SPGreen!70,rounded corners=.8mm,fill=SPGreen!5,
 font=\scriptsize,minimum width=3.3cm,minimum height=.58cm] at (3,-1.25)
 {$\mathcal B=\{p_1,p_4\}$: independiente, no necesariamente máximo};
\end{tikzpicture}
}
\caption{Las propuestas físicas inducen un grafo de conflictos; DMIS selecciona
un conjunto independiente maximal, no un conjunto máximo.}
\label{fig:n4-conflict-graph}
\viuownsource
\end{figure}
\vspace{-3mm}

CF conserva un orden global como ablación; DMIS lo sustituye por prioridades
entre vecinos de conflicto~\parencite{luby1986MIS}. Cada carga usa
$z_k=(Q_k,v_k^{\mathrm{reg}})$, un registro lógico del simulador. Un
\emph{keeper} determinista recibe propuestas y propaga cada actualización por
un árbol vecinal. Los registros viven en el mismo proceso: DMIS+TX elimina el
orden global, no toda infraestructura lógica.

\begin{figure}[H]
\centering
\scalebox{.72}{\begin{tikzpicture}[x=1cm,y=1cm,>=Latex,
 stage/.style={draw=SPGrid,rounded corners=1mm,fill=SPFloor,text width=1.75cm,
  minimum height=1.05cm,align=center,font=\scriptsize},
 arr/.style={->,thick,draw=VIUOrange}]
\node[stage,draw=SPBlue!65] (read) at (-4.7,.65) {\textbf{READ}\\$z_r,z_s$};
\node[stage,draw=SPPurple!65] (prop) at (-2.35,.65) {\textbf{PROPOSE}\\$p,\Delta\Psi,v_r,v_s$};
\node[stage,draw=SPGreen!65] (mis) at (0,.65) {\textbf{DMIS}\\vecinos de conflicto};
\node[stage,draw=VIUOrange!75] (prep) at (2.35,.65) {\textbf{PREPARE}\\validar y bloquear};
\node[stage,draw=SPGreen!70,fill=SPGreen!6] (commit) at (4.7,1.25) {\textbf{COMMIT}\\escribir todo};
\node[stage,draw=SPRed!70,fill=SPRed!5] (abort) at (4.7,-.35) {\textbf{ABORT}\\liberar reservas};
\draw[arr] (read)--(prop); \draw[arr] (prop)--(mis); \draw[arr] (mis)--(prep);
\draw[arr] (prep)--(commit); \draw[->,thick,draw=SPRed] (prep)--(abort);
\draw[->,thin,draw=SPBlue] (commit.north) |- ++(0,.55) -| (read.north);
\node[font=\tiny,text=black!55] at (0,-.72)
 {sin orden global de confirmación};
\end{tikzpicture}
}
\caption{Geo-QPG-DMIS+TX separa lectura, propuesta, arbitraje vecinal y
\emph{commit} versionado.}
\label{fig:n4-dmis-architecture}
\viuownsource
\end{figure}
\vspace{-3mm}

\begin{tcolorbox}[colback=SPFloor,colframe=black!18,boxrule=.45pt,arc=.8mm,
 left=2mm,right=2mm,top=.6mm,bottom=.6mm]
\footnotesize La ejecución sigue cinco operaciones: \texttt{READ} lee
estado y versiones; \texttt{PROPOSE} publica la mejora; \texttt{DMIS} elige
ganadores vecinales; \texttt{PREPARE/COMMIT} escribe todo si las versiones siguen
vigentes; y \texttt{ABORT} no escribe.
\end{tcolorbox}

El \emph{commit} escribe todo o nada solo con entrega fiable y proponente activo.
Pérdida, retardo, partición y fallos bizantinos quedan fuera
\parencite{grayLamport2006Commit}. F-II y F-III recuperan la coalición entera al
cierre.

\clearpage

% --------------------------------------------------------------------------
% Página 31 - familia poblacional
% --------------------------------------------------------------------------
\leveltag{RESULTADOS SP1 · N4 · F-II · JUEGOS POBLACIONALES}
\section*{F-II: intención continua y cierre atómico}
\enlargethispage{1\baselineskip}

F-II sustituye la acción indivisible por $x_i\in\Delta_i$, con
$x_{ik}\in[0,1]$ y $\sum_{k=0}^{K}x_{ik}=1$. La capacidad relajada
$Q_k(x)=\sum_i c_i^{\mathrm{pay}}x_{ik}$ no identifica todavía qué robots
acudirán. E7 (véase más adelante en este nivel) fija una sola relajación para
que las cuatro dinámicas vean el mismo campo. Sea
$r_k(x)=[1-Q_k(x)/m_k]_+$ y
$\bar d_{ik}=d_{ik}/\max_{r,s}d_{rs}$:
\begin{equation}
\label{eq:n4-smooth-potential}
\widetilde\Phi(x)=
-\frac{\alpha}{2}\sum_k r_k(x)^2
-\frac{\beta}{N}\sum_{i,k}\bar d_{ik}x_{ik},
\qquad
F_{ik}=\frac{\partial\widetilde\Phi}{\partial x_{ik}}.
\end{equation}
La penalización cuadrática de déficit (\emph{squared hinge}) es de clase $C^1$ y
su gradiente se anula al cubrir la cuota. La normalización hace explícito el
compromiso $\alpha/\beta$; no reproduce el
orden lexicográfico de F-I y se trata como una penalización continua.

Replicator, Smith, Brown--von Neumann--Nash (BNN) y Logit cambian la regla de
revisión, no el campo~\parencite{sandholm2010population,quijano2017population}:
\begin{equation}
\label{eq:n4-population-dynamics}
\begin{aligned}
\dot x_{ik}^{\mathrm{Rep}}&=x_{ik}(F_{ik}-\bar F_i),
&\bar F_i&=\sum_qx_{iq}F_{iq},\\[-1mm]
\dot x_{ik}^{\mathrm{Smith}}&=\sum_qx_{iq}[F_{ik}-F_{iq}]_+
-x_{ik}\sum_q[F_{iq}-F_{ik}]_+,\\[-1mm]
\dot x_{ik}^{\mathrm{BNN}}&=[F_{ik}-\bar F_i]_+
-x_{ik}\sum_q[F_{iq}-\bar F_i]_+,\\[-1mm]
\dot x_{ik}^{\mathrm{Logit}}&=
\frac{\exp(F_{ik}/\tau)}{\sum_q\exp(F_{iq}/\tau)}-x_{ik}.
\end{aligned}
\end{equation}

\begin{figure}[H]
\centering
\begin{tikzpicture}[x=1cm,y=1cm,>=Latex,
 box/.style={draw=SPGrid,rounded corners=1mm,fill=SPFloor,
  minimum height=1.55cm,align=center,font=\scriptsize},
 robot/.style={circle,draw=SPBlue,fill=white,minimum size=5.5mm,font=\tiny},
 load/.style={draw=VIUOrange,rounded corners=.6mm,fill=VIUOrange!5,
  minimum width=7mm,minimum height=5.5mm,font=\tiny},
 arr/.style={->,very thick,draw=VIUOrange}]
\node[box,text width=2.25cm,draw=SPBlue!65] (world) at (-4.55,.55)
 {\textbf{mundo físico}\\[1mm]
  \tikz[baseline=-.6ex]{\node[robot] {$r_1$};}\quad
  \tikz[baseline=-.6ex]{\node[robot] {$r_2$};}\quad
  \tikz[baseline=-.6ex]{\node[robot] {$r_3$};}\\[-.2mm]
  \tikz[baseline=-.5ex]{\node[load] {$k_1$};}\qquad
  \tikz[baseline=-.5ex]{\node[load] {$k_2$};}};

\node[box,text width=2.65cm,draw=SPGreen!70] (simplex) at (-1.55,.55)
 {\textbf{intenciones $x_i$}\\[-.2mm]
  \begin{tikzpicture}[x=.16cm,y=.16cm]
   \fill[SPBlue!65] (0,4) rectangle (7,5.5); \fill[SPGreen!65] (7,4) rectangle (12,5.5);
   \fill[SPBlue!65] (0,1.5) rectangle (3,3); \fill[SPGreen!65] (3,1.5) rectangle (12,3);
   \fill[SPBlue!65] (0,-1) rectangle (6,-.5); \fill[SPGreen!65] (6,-1) rectangle (12,-.5);
   \draw (0,-1) rectangle (12,5.5);
  \end{tikzpicture}};

\node[box,text width=2.45cm,draw=SPPurple!70] (flow) at (1.45,.55)
 {\textbf{mismo campo $F$}\\Replicator · Smith\\BNN · Logit\\$x(t)\rightarrow x^\star$};

\node[box,text width=2.45cm,draw=VIUOrange!80] (closure) at (4.45,.55)
 {\textbf{cierre común}\\$a^\star=\mathcal R(x^\star)$\\[1mm]
  \tikz[baseline=-.6ex]{\node[robot,draw=SPGreen] {$r_1$};}
  $\rightarrow$
  \tikz[baseline=-.5ex]{\node[load,draw=SPGreen] {$k_1$};}\\coalición entera};

\draw[arr] (world)--(simplex); \draw[arr] (simplex)--(flow); \draw[arr] (flow)--(closure);
\node[draw=SPRed!65,rounded corners=.8mm,fill=SPRed!3,
 text width=8.7cm,minimum height=.62cm,align=center,font=\scriptsize] at (1.45,-1.10)
 {$x_{ik}$ es masa/probabilidad; \textbf{frontera de evidencia:} converger en el símplex no implica calidad después de $\mathcal R$};
\end{tikzpicture}

\caption{F-II mantiene un estado divisible durante la evolución y requiere un
cierre común antes de producir robots enteros.}
\label{fig:n4-population-closure}
\viuownsource
\end{figure}

La ejecución distribuida mantiene una estimación mediante consenso de media
dinámica (DAC) de $Q_k(x)$ y registra su
error. Replicator conserva soporte; Logit explora mediante $\tau$. La siguiente
página fija el algoritmo digital y el cierre común empleado en E7.

\clearpage

% --------------------------------------------------------------------------
% Página 32 - familia Nash restringida
% --------------------------------------------------------------------------
\leveltag{RESULTADOS SP1 · N4 · F-III · PRIMAL--DUAL Y EQUILIBRIO RESTRINGIDO}
\section*{F-III: el déficit se trata como restricción compartida}
\enlargethispage{1\baselineskip}
\begingroup
\setlength{\abovedisplayskip}{5pt}
\setlength{\belowdisplayskip}{5pt}
\setlength{\abovedisplayshortskip}{3pt}
\setlength{\belowdisplayshortskip}{3pt}

F-III también usa $x$, pero convierte cada cuota en una restricción acoplada
\begin{equation}
\label{eq:n4-shared-constraint}
g_k(x)=1-\frac{Q_k(x)}{m_k}\le0,
\quad
J_i(x_i)=\frac1N\sum_k\bar d_{ik}x_{ik}
+\frac{\epsilon}{2}\lVert x_i\rVert^2,
\quad
\mathcal L=\sum_iJ_i+\sum_k\lambda_kg_k.
\end{equation}
El juego queda definido por los $N$ problemas
\begin{equation}
\label{eq:n4-vgne-player-problem}
\min_{x_i\in\Delta_i}J_i(x_i)
\quad\text{s.a.}\quad g(x)\leq0,qquad i=1,\ldots,N,
\end{equation}
con $\mathcal X=\prod_i\Delta_i$,
$\mathbf F(x)=\operatorname{col}_i(\nabla_{x_i}J_i)$ y Lagrangiano individual
$\mathcal L_i=J_i+\lambda^\top g$. Con restricciones compartidas y potencial
convexo, el vGNE de multiplicador común cumple las KKT del programa asociado.
$\lambda_k\ge0$ actúa como precio de capacidad: si falta cobertura, el término
$(c_i^{\mathrm{pay}}/m_k)\lambda_k$ atrae masa hacia $k$. La regularización
$\epsilon=0{,}02$ hace estrictamente convexo el problema continuo.

La referencia central resuelve el programa convexo. La versión distribuida
discretiza el mecanismo primal--dual:
\begin{equation}
\label{eq:n4-primal-dual-candidate}
\dot x_i=\Pi_{T_{\Delta_i}(x_i)}[-\nabla_{x_i}\mathcal L_i],
\qquad
\dot\lambda_k=\Pi_{T_{\mathbb R_+}(\lambda_k)}[g_k(x)].
\end{equation}
Cada robot mantiene $\lambda_{ik}$ y una estimación
$\widehat g_{ik}$, con un término explícito de consenso:
\begin{equation}
\label{eq:n4-distributed-dual-candidate}
\dot\lambda_{ik}=\Pi_{T_{\mathbb R_+}(\lambda_{ik})}
\left[\widehat g_{ik}-\gamma\sum_{j\in\mathcal N_i}
(\lambda_{ik}-\lambda_{jk})\right].
\end{equation}

\begin{figure}[H]
\centering
\vspace{-3mm}
\scalebox{.58}{\begin{tikzpicture}[x=1cm,y=1cm,>=Latex,
 market/.style={draw=VIUOrange!80,rounded corners=1mm,fill=VIUOrange!5,
  minimum width=1.9cm,minimum height=1.15cm,align=center,font=\scriptsize},
 robot/.style={circle,draw=SPBlue,fill=white,minimum size=6.2mm,font=\scriptsize},
 box/.style={draw=SPGrid,rounded corners=1mm,fill=SPFloor,
  minimum height=1.25cm,align=center,font=\scriptsize},
 arr/.style={->,thick,draw=VIUOrange}]
\node[market] (k1) at (-3.9,1.35) {carga $k_1$\\$g_1>0\Rightarrow\lambda_1\uparrow$};
\node[market] (k2) at (0,1.35) {carga $k_2$\\$g_2\simeq0$};
\node[market] (k3) at (3.9,1.35) {carga $k_3$\\$g_3<0$};
\node[robot] (r1) at (-3.1,-.15) {$i$};
\node[robot] (r2) at (-1.05,-.15) {$j$};
\node[robot] (r3) at (1.05,-.15) {$\ell$};
\node[robot] (r4) at (3.1,-.15) {$q$};
\draw[dashed,draw=SPPurple,thick] (r1)--(r2)--(r3)--(r4);
\node[font=\tiny,text=SPPurple] at (0,-.72)
 {consenso/DAC entre copias locales $\lambda_{ik}$ y estimaciones $\widehat g_{ik}$};
\foreach \r/\k in {r1/k1,r2/k1,r2/k2,r3/k2,r3/k3,r4/k3}{\draw[arr] (\r)--(\k);}
\node[box,text width=3.05cm,draw=SPPurple!70] (pd) at (-2.1,-1.55)
 {\textbf{mecanismo primal--dual}\\$-d_{ik}+c_i^{\rm pay}\lambda_{ik}$};
\node[box,text width=2.15cm,draw=SPGreen!70] (vgne) at (1.0,-1.55)
 {\textbf{objetivo}\\vGNE continuo $(x^\star,\lambda^\star)$};
\node[box,text width=2.15cm,draw=VIUOrange!80] (close) at (4.0,-1.55)
 {\textbf{salida física}\\$\mathcal R(x^\star)=a^\star$};
\draw[arr] (pd)--(vgne); \draw[arr] (vgne)--(close);
\end{tikzpicture}
}
\caption{Interpretación de F-III: las cargas deficitarias elevan su precio,
las copias duales se coordinan y el estado continuo necesita cierre atómico.}
\label{fig:n4-primal-dual}
\viuownsource
\end{figure}
\vspace{-3mm}

El vGNE con multiplicador común satisface:
\begin{equation}
\label{eq:n4-vgne-kkt}
\begin{gathered}
0\in\mathbf F(x^\star)+\nabla g(x^\star)^\top\lambda^\star
+N_{\mathcal X}(x^\star),\\[-1mm]
g(x^\star)\le0,\quad\lambda^\star\ge0,\quad
{\lambda^\star}^{\!\top}g(x^\star)=0.
\end{gathered}
\end{equation}

El criterio numérico reúne las cinco condiciones. La literatura motiva el
mecanismo~\parencite{cherukuri2016primalDual,yi2019operatorGNE}; no se afirma
convergencia global de la discretización evaluada.
\endgroup

\clearpage

% --------------------------------------------------------------------------
% Página 33 - algoritmos continuos y cierre común
% --------------------------------------------------------------------------
\leveltag{RESULTADOS SP1 · N4 · F-II--F-III · ALGORITMOS Y CIERRE}
\section*{Algoritmos continuos y cierre entero común}

La iteración es digital y usa rondas simuladas sobre el grafo estático de N3;
no se presenta como asíncrona. Los dos mecanismos terminan en el mismo
$\mathcal R$ para que la comparación no dependa de un redondeo distinto por
familia.

\begin{tcolorbox}[colback=SPBlue!3,colframe=SPBlue!65,boxrule=.6pt,arc=1mm,
 left=2mm,right=2mm,top=1.4mm,bottom=1.4mm]
\textbf{Algoritmo 4 · F-II distribuido con dinámica $M$}
\begin{algorithmic}[1]
\Require mundo, grafo $G$, $M\in\{\mathrm{Rep,Smith,BNN,Logit}\}$, $x^0$, $T$
\State cada robot inicia $z_{ik}^0=(c_i^{\mathrm{pay}}/m_k)x_{ik}^0$ para todo $k$
\For{$t=0,\ldots,T-1$}
  \State intercambiar $z_i^t$ con vecinos y actualizar DAC
  \State calcular $\widehat Q_{ik}^t$ y el mismo fitness $F_{ik}(x^t)$
  \State obtener $v_i^t=M(x_i^t,F_i^t)$ y proyectar $x_i^{t+1}$ en $\Delta_i$
  \State corregir $z_i^{t+1}$ con el cambio de contribución local
  \If{residual de punto fijo y error DAC cruzan sus umbrales} \State \textbf{salir} \EndIf
\EndFor
\State \Return $x^\star$, historial y $a^\star=\mathcal R(x^\star)$
\end{algorithmic}
\end{tcolorbox}

\begin{tcolorbox}[colback=SPGreen!3,colframe=SPGreen!65,boxrule=.6pt,arc=1mm,
 left=2mm,right=2mm,top=1.4mm,bottom=1.4mm]
\textbf{Algoritmo 5 · F-III primal--dual distribuido}
\begin{algorithmic}[1]
\Require mundo, grafo $G$, $x^0$, copias $\lambda_i^0=0$, ganancias y $T$
\For{$t=0,\ldots,T-1$}
  \State mezclar por vecinos las copias $\lambda_i^t$ y el estimador de $g_k(x^t)$
  \State predictor: proyectar el paso primal y actualizar $\lambda_i\ge0$
  \State corrector: reevaluar gradiente, DAC y copias duales
  \State calcular residuos primal, dual, complementariedad, estación y consenso
  \If{la norma conjunta es menor que $5\times10^{-4}$} \State \textbf{salir} \EndIf
\EndFor
\State \Return $(x^\star,\lambda^\star)$, historial y $a^\star=\mathcal R(x^\star)$
\end{algorithmic}
\end{tcolorbox}

\begin{tcolorbox}[colback=SPFloor,colframe=VIUOrange!80,boxrule=.6pt,arc=1mm,
 left=2mm,right=2mm,top=1.3mm,bottom=1.3mm]
\small
\textbf{Cierre común $\mathcal R$.} Es un postproceso central del observador,
común a F-II y F-III, no una fase distribuida. Ordena para cada robot las cargas por
$x_{ik}$, conserva asignaciones compatibles y aplica cadenas aumentantes
deterministas hasta cubrir cuotas o agotar límites. Devuelve factibilidad,
distancia, exceso y operaciones. Es una heurística acotada; el MILP sigue siendo
el único certificado entero global. Por ello estas ramas comparan una intención
distribuida seguida del mismo cierre atómico; no son implementaciones
distribuidas extremo a extremo.
\end{tcolorbox}

\clearpage

% --------------------------------------------------------------------------
% Página 34 - familia estocástica y mismo mundo
% --------------------------------------------------------------------------
\leveltag{RESULTADOS SP1 · N4 · F-IV · DECISIÓN RECURRENTE}
\section*{F-IV: reclutar de nuevo cuando cambia el mundo}
\enlargethispage{2\baselineskip}

Las tres primeras familias resuelven una instantánea. Cuando llegan trabajos,
cambia la disponibilidad, cae un robot o termina un servicio, el reclutamiento
pasa a ser secuencial. E8 usa el modelo mínimo que permite comprobar esa idea
sin invadir la dinámica física de SP2:
\begin{equation}
\label{eq:n4-markov-formulation}
\begin{aligned}
s_t&=(a_{t-1},\mathcal J_t,\mathcal A_t,e_t),
\qquad u_t\in\{0,\ldots,K\}^{N},\\
s_{t+1}&=T(s_t,u_t,e_{t+1}),\\
C_t&=w_D\bar D_t+w_d\bar J_t+w_s
\lVert u_t-a_{t-1}\rVert_0,
\qquad V^\pi(s_0)=\sum_{t=0}^{T-1}\gamma^tC_t .
\end{aligned}
\end{equation}
Los eventos $e_t$ son llegada, finalización o fallo y se aplican antes de cada
decisión. Se comparan Geo-QPG desde cero, Geo-QPG activado por evento e
inicializado desde la solución anterior, y un programa dinámico central que
conoce la secuencia completa. $T$ es una actualización determinista condicionada
por el evento observado; no se modela probabilísticamente ni se estima un kernel.
Franci y
Grammatico estudian GNE estocástico bajo información parcial
\parencite{franci2022stochasticGNE}; sus garantías no se transfieren al estudio
temporal evaluado aquí.

\begin{figure}[H]
\centering
\begin{tikzpicture}[x=1cm,y=1cm,>=Latex,
 panel/.style={draw=SPGrid,rounded corners=1mm,fill=SPFloor,
  minimum width=5.15cm,minimum height=2.55cm,align=center,font=\scriptsize},
 robot/.style={circle,draw=SPBlue,fill=white,minimum size=5mm,font=\tiny},
 load/.style={draw=VIUOrange,rounded corners=.6mm,fill=VIUOrange!5,
  minimum width=6.5mm,minimum height=5mm,font=\tiny},
 arr/.style={->,thick,draw=VIUOrange}]
\foreach \x/\y/\name/\title in {-2.8/1.55/f1/{F-I · entero},2.8/1.55/f2/{F-II · intención},-2.8/-1.35/f3/{F-III · precio},2.8/-1.35/f4/{F-IV · evento}}{
  \node[panel] (\name) at (\x,\y) {};
  \node[font=\scriptsize\bfseries,anchor=north west] at ($(\name.north west)+(.18,-.15)$) {\title};
  \node[robot] (\name-r1) at ($(\name.center)+(-1.45,.35)$) {$1$};
  \node[robot] (\name-r2) at ($(\name.center)+(-.15,.65)$) {$2$};
  \node[robot] (\name-r3) at ($(\name.center)+(.15,-.45)$) {$3$};
  \node[load] (\name-k1) at ($(\name.center)+(1.25,.35)$) {$k_1$};
  \node[load] (\name-k2) at ($(\name.center)+(1.65,-.45)$) {$k_2$};
}
\draw[arr] (f1-r1)--(f1-k1); \draw[arr] (f1-r2)--(f1-k1); \draw[arr] (f1-r3)--(f1-k2);
\node[font=\tiny,text=SPBlue] at ($(f1.center)+(0,-.92)$) {$a_i$ permanece indivisible};

\draw[arr,dashed] (f2-r1)--node[pos=.25,above,sloped,fill=white,inner sep=.2mm,font=\tiny] {$0{,}7$}(f2-k1);
\draw[arr,dashed] (f2-r1)--node[pos=.27,below,sloped,fill=white,inner sep=.2mm,font=\tiny] {$0{,}3$}(f2-k2);
\draw[arr,dashed] (f2-r2)--(f2-k1); \draw[arr,dashed] (f2-r3)--(f2-k2);
\node[font=\tiny,text=SPGreen!70!black] at ($(f2.center)+(0,-.92)$) {$x^\star\xrightarrow{\mathcal R}a$};

\node[font=\tiny,text=SPPurple] at ($(f3-k1.north)+(0,.18)$) {$\lambda_1\uparrow$};
\node[font=\tiny,text=SPPurple] at ($(f3-k2.north)+(0,.18)$) {$\lambda_2$};
\draw[dashed,draw=SPPurple] (f3-r1)--(f3-r2)--(f3-r3);
\draw[arr] (f3-r1)--(f3-k1); \draw[arr] (f3-r2)--(f3-k1); \draw[arr] (f3-r3)--(f3-k2);
\node[font=\tiny,text=SPPurple] at ($(f3.center)+(0,-.92)$) {primal--dual $\rightarrow$ vGNE $\rightarrow\mathcal R$};

\node[robot,draw=SPRed] (failed) at (f4-r2) {$\times$};
\node[load,draw=SPRed,dashed] (new) at ($(f4.center)+(.70,-.45)$) {$k_3$};
\draw[arr] (f4-r1)--(f4-k1); \draw[arr] (f4-r3)--(new);
\node[font=\tiny,text=SPRed] at ($(f4.center)+(0,-1.02)$) {$t_0\rightarrow$ fallo/llegada $\rightarrow t_1$: nueva política};
\end{tikzpicture}

\caption{El mismo mundo físico visto por las cuatro familias. F-I mantiene
robots enteros; F-II y F-III requieren cierre; F-IV añade evolución temporal.}
\label{fig:n4-same-world}
\viuownsource
\end{figure}

\begin{tcolorbox}[colback=SPFloor,colframe=VIUOrange!80,boxrule=.65pt,arc=1mm,
 left=2.2mm,right=2.2mm,top=1.6mm,bottom=1.6mm]
\small
E8 se modela como un juego potencial repetido condicionado por eventos exógenos.
No se estima un kernel $P$ ni se ha probado una propiedad de potencial de Markov.
El oráculo del objetivo escalar con información futura es central, de horizonte
finito y viable solo para el tamaño pequeño enumerado.
\end{tcolorbox}

La comparación se hace sobre salidas enteras: F-I entrega $a$ directamente,
F-II y F-III pasan por $\mathcal R$, y F-IV produce una secuencia de perfiles.
E8 no comparte la instantánea estática de E7 y, por ello, no entra en su ranking.

\clearpage

% --------------------------------------------------------------------------
% Página 35 - información, cadena formal y programa experimental
% --------------------------------------------------------------------------
\leveltag{RESULTADOS SP1 · N4 · INFORMACIÓN Y PROGRAMA EXPERIMENTAL}
\section*{Información intercambiada y diseño experimental}
\enlargethispage{2\baselineskip}

La red transporta estados distintos en cada familia. F-I intercambia $Q_k$, la
versión y la propuesta sobre las cargas afectadas. F-II estima $Q_k(x)$ y el
fitness mediante DAC; F-III añade copias de los multiplicadores y consenso
dual. F-IV transmite eventos y disponibilidad, pero su oráculo temporal sigue
siendo central. Estos mecanismos determinan qué información está disponible;
no cambian por sí solos el concepto de solución.

En F-I ya se establecieron la suficiencia de los agregados, el potencial exacto
unilateral, la terminación en una vecindad finita y la composición de propuestas
no conflictivas. La confirmación atómica supone entrega fiable. Por ello,
la cadena formal no cubre optimalidad global, Nash fuerte, fallos de red ni
transporte físico.

\begin{table}[H]
\centering
\caption{Diseño de los experimentos de N4 y variables analizadas.}
\label{tab:n4-evidence-programme}
\fontsize{7.7}{8.5}\selectfont
\setlength{\tabcolsep}{1.8pt}
\renewcommand{\arraystretch}{1.08}
\begin{tabularx}{\textwidth}{>{\raggedright\arraybackslash}p{.62cm}
 >{\raggedright\arraybackslash}p{3.15cm}Y
 >{\raggedright\arraybackslash}p{2.25cm}}
\toprule
Estudio & Contraste & Variable principal & Resultado gráfico \\
\midrule
E4 & BR--ASR--LLL; después BR--2BR--C3 & brecha, factibilidad y tráfico &
Fig.~\ref{fig:n4-family-results} \\
E4 & CF frente a DMIS+TX & efecto del orden global de confirmación &
Fig.~\ref{fig:n4-architecture-results} \\
E9 & búsqueda conectada hasta $h=3$ & primer orden conectado $h_{\mathrm c}^\star$ &
Figs.~\ref{fig:n4-hstar-distribution}--\ref{fig:n4-fi-trace} \\
E5 & DPOP frente a MILP en tamaño enumerable & exactitud y tamaño de tabla &
Fig.~\ref{fig:n4-dpop-results} \\
E6 & BR frente a DMIS+TX, $16\le N\le64$ & tiempo y bytes por agente &
Fig.~\ref{fig:n4-scaling-results} \\
E7 & Replicator, Smith, BNN y Logit con el mismo cierre & factibilidad y brecha &
Fig.~\ref{fig:n4-fii-continuous-atomic} \\
E7 & primal--dual frente a Smith con el mismo cierre & KKT, factibilidad y tráfico &
Fig.~\ref{fig:n4-fiii-residuals} \\
E8 & políticas reactivas y oráculo temporal & coste, servicio y cambios &
Fig.~\ref{fig:n4-fiv-dynamic} \\
\bottomrule
\end{tabularx}
\viuownsource
\end{table}

{\small El mundo--semilla es la unidad independiente; los fallos y los límites
temporales permanecen en el denominador. E4--E6 y E9 estudian F-I, E7 compara
las tres familias estáticas con salida entera y E8 analiza la decisión temporal.
El canal degradado queda fuera del dominio experimental de este capítulo.\par}

\clearpage

% --------------------------------------------------------------------------
% Página 36 - resultado de orden y revisión
% --------------------------------------------------------------------------
\leveltag{RESULTADOS SP1 · N4 · E4 · ORDEN Y REVISIÓN}
\section*{E4: la mejora provino de aumentar el orden local}

E4 usa un banco confirmatorio disjunto de 1\,200 mundos, generado con los mismos
factores de N3 pero con una corriente de semillas propia. Los tres baselines de
N3 y las variantes N4 reciben dentro de cada mundo las mismas capacidades,
demandas y geometría; en total se obtuvieron \NFourRows{} ejecuciones. Así se
separan tres intervenciones: el orden $h$, la regla de revisión y la arquitectura
DMIS+TX (Figura~\ref{fig:n4-family-results}).

\begin{figure}[H]
\centering
\includegraphics[width=.90\linewidth]{assets/figures/n4-v2/n4_family_comparison.pdf}
\caption{E4: calidad y coste al ampliar la vecindad local. Los intervalos se
calculan sobre mundos con óptimo MILP certificado.}
\label{fig:n4-family-results}
\viuownsource
\end{figure}

Al ampliar el orden local, la brecha (\emph{gap}) mediana bajó de
\NFourGapPctQpgU\,\% ($h=1$) a \NFourGapPctQpgP\,\% ($h=2$) y
\NFourGapPctCThree\,\% ($h=3$). C3$-$BR fue \NFourCThreeGapDiffPct{} puntos
(IC 95\,\%: \NFourCThreeGapDiffLowPct{} a \NFourCThreeGapDiffHighPct{}), elevó
la factibilidad \NFourCThreeRiskPct{} puntos y costó \NFourBytesCThree{} bytes
por agente y \NFourRuntimeCThree{} ms.

Con $h=1$, Geo-ASR y Geo-LLL empeoraron la brecha de BR en
\NFourSmithGapDiffPct{} y \NFourLllGapDiffPct{} puntos. En estos mundos, cambiar
la regla de revisión no compensó la falta de coordinación entre dos o tres
robots. E9 examina a continuación cuál fue el primer orden de escape de cada
terminal BR.

\begin{table}[H]
\centering
\caption{Mapa operativo de la familia F-I, con el alcance de cada variante.}
\label{tab:n4-fi-methods}
\fontsize{7.2}{7.8}\selectfont
\setlength{\tabcolsep}{2.1pt}
\renewcommand{\arraystretch}{0.92}
\begin{tabularx}{\textwidth}{>{\raggedright\arraybackslash}p{1.25cm}c
 >{\raggedright\arraybackslash}p{1.55cm}
 >{\raggedright\arraybackslash}p{1.35cm}
 >{\raggedright\arraybackslash}p{2.35cm}c
 >{\raggedright\arraybackslash}p{1.55cm}Y}
\toprule
Método & $h$ & Revisión & Scheduler & Información & Atómico & Terminal & Coste principal \\
\midrule
BR & 1 & mejor respuesta & secuencial & cargas afectadas & sí & 1-local & calidad limitada \\
2BR & 2 & mejor par & secuencial & cargas + pareja & sí & 2-local & más candidatos \\
C3 & 3 & mejor triple & secuencial & cargas + triple & sí & 3-local & cómputo y bytes \\
Geo-ASR & 1 & proporcional & secuencial & cargas afectadas & sí & 1-local & muestreo estocástico \\
Geo-LLL & 1 & Gibbs + pulido BR & secuencial & cargas afectadas & sí & 1-local si no trunca & explora y después desciende \\
CF & 1--3 & candidatos QPG & orden global & lote completo & sí & $h$-local & coordinación global \\
DMIS+TX & 1--3 & candidatos QPG & vecinal & conflictos + versiones & sí & $h$-local & tráfico adicional \\
\bottomrule
\end{tabularx}
\viuownsource
\end{table}

\clearpage

% --------------------------------------------------------------------------
% E9 - distribución del orden mínimo
% --------------------------------------------------------------------------
\leveltag{RESULTADOS SP1 · N4 · E9 · ORDEN MÍNIMO DE COORDINACIÓN}
\section*{E9: en los mundos evaluados, el primer escape conectado fue bilateral}
\enlargethispage{2\baselineskip}

Después de ejecutar BR desde el perfil ocioso, se reconstruyeron los
\NFourHStarWorlds{} mundos y se buscó el primer grupo conectado que redujese
$(D_4,J_4)$. Por construcción, $P(h_{\mathrm c}^\star=1)=0$ comprueba que BR agotó su
vecindad. La búsqueda terminó en tres; ``$>3$'' solo indica que no encontró una
mejora dentro de ese corte.

\begin{figure}[H]
\centering
\includegraphics[width=.79\linewidth]{assets/figures/n4-v4/n4_hstar_distribution.pdf}
\caption{E9: distribución de $h_{\mathrm c}^\star$, el primer orden conectado
que mejora el terminal BR.
Cada barra contiene 150 mundos en el panel de presión--heterogeneidad y 240 por
geometría.}
\label{fig:n4-hstar-distribution}
\viuownsource
\end{figure}

El primer escape fue bilateral en \NFourHStarTwoPct\,\% de los mundos y
trilateral en \NFourHStarThreePct\,\%; en el
\NFourHStarGtThreePct\,\% restante no apareció una mejora hasta orden tres. En
este banco, el bloqueo que dejó BR se resolvió casi siempre mediante un cambio
de dos robots. La búsqueda comienza en el terminal obtenido desde reposo y se
detiene en $h=3$, de modo que la última categoría no certifica optimalidad.

\begin{figure}[H]
\centering
\includegraphics[width=.76\linewidth]{assets/figures/n4-v4/n4_hstar_ratio_sensitivity.pdf}
\caption{Sensibilidad independiente de E9: proporción sin mejora conectada hasta
$h=3$ al variar $N/K$, presión y heterogeneidad. Este bloque no se mezcla con
los 1\,200 mundos confirmatorios.}
\label{fig:n4-hstar-ratio}
\viuownsource
\end{figure}

La fracción $h_{\mathrm c}^\star\ge3$ pasó del \NFourHStarHardLowCvPct\,\% con $CV=0$ al
\NFourHStarHardHighCvPct\,\% con $CV=1$; la categoría $>3$ cayó del
\NFourHStarGtThreeLowRatioPct\,\% con $N/K=2$ al
\NFourHStarGtThreeHighRatioPct\,\% con $N/K=4$. Son tendencias del barrido, no
monotonía ni una ley asintótica.

\clearpage

% --------------------------------------------------------------------------
% E9 - enlace entre barrera, mejora y ejecución
% --------------------------------------------------------------------------
\leveltag{RESULTADOS SP1 · N4 · E9 · MECANISMO Y TRAZA}
\section*{E9: orden de escape, coste y traza de ejecución}
\enlargethispage{4\baselineskip}

El orden de escape también anticipó qué refinamiento podía mejorar el terminal.
En los mundos con escape bilateral, 2BR redujo la brecha mediana en
\NFourHStarPairGainPct{} puntos respecto de BR. En los cuatro casos con
$h_{\mathrm c}^\star=3$, C3 aportó otros \NFourHStarTripleGainPct{} puntos respecto de 2BR.
Cuando no apareció una mejora hasta tres, ninguno de los dos refinamientos
redujo el perfil BR.

\begin{figure}[H]
\centering
\includegraphics[width=.84\linewidth]{assets/figures/n4-v4/n4_hstar_gap_link.pdf}
\caption{Relación entre el orden de escape y la reducción observada de la brecha. La
línea negra es la mediana; los puntos son mundos y la anchura resume densidad.}
\label{fig:n4-hstar-gap-link}
\viuownsource
\end{figure}

\begin{figure}[H]
\centering
\includegraphics[width=.76\linewidth]{assets/figures/n4-v4/n4_fi_representative_trace.pdf}
\caption{Traza de F-I: déficit y distancia después de cada confirmación,
orden de las desviaciones aceptadas y tráfico total. La distancia puede crecer
durante la fase de cuota porque el orden lexicográfico prioriza eliminar déficit.}
\label{fig:n4-fi-trace}
\viuownsource
\end{figure}

C3 alcanzó una distancia menor mediante movimientos de orden dos y tres, pero
exploró más candidatos y acumuló más bytes; la traza conserva cada cambio
aceptado. El orden exacto de los 40 casos truncados sigue sin conocerse. Tampoco
se desprende de E9 que el terminal de C3 sea globalmente óptimo.

\clearpage

% --------------------------------------------------------------------------
% Página 37 - arquitectura y exactitud
% --------------------------------------------------------------------------
\leveltag{RESULTADOS SP1 · N4 · E4--E5 · ARQUITECTURA Y EXACTITUD}
\section*{Arquitectura vecinal y exactitud en pequeño}
\enlargethispage{1\baselineskip}

Aquí se responden dos preguntas distintas. Primero se mantienen los candidatos y
se retira el orden global de confirmación. Después se contrasta, en tamaños
enumerables, un método exacto distribuido sobre pseudoárbol
(DPOP)\parencite{petcuFaltings2005DPOP} con el MILP.

\begin{figure}[H]
\centering
\includegraphics[width=.68\linewidth]{assets/figures/n4-v2/n4_architecture_ablation.pdf}
\caption{E4: ablación arquitectónica CF frente a DMIS+TX sobre los mismos
mundos y candidatos.}
\label{fig:n4-architecture-results}
\viuownsource
\end{figure}

En el canal nominal, retirar el orden global conservó la factibilidad y la
brecha: las diferencias pareadas fueron \NFourDRiskPct{} puntos y
\NFourDGapDiffPct{}, respectivamente. DMIS+TX añadió \NFourDByteDiff{} bytes por
agente (IC 95\,\%: \NFourDByteDiffLow{} a \NFourDByteDiffHigh{}). Estas cifras
son contadores de campos por salto bajo el canal nominal. No se midieron rondas,
identidad final o latencia física, ni se inyectaron fallos.

\begin{figure}[H]
\centering
\includegraphics[width=.55\linewidth]{assets/figures/n4-v2/n4_dpop_exactness.pdf}
\caption{E5: exactitud de DPOP y crecimiento de su tabla UTIL.}
\label{fig:n4-dpop-results}
\viuownsource
\end{figure}

DPOP se evaluó con $N\in\{6,8\}$, $K=2$, cinco geometrías,
$CV\in\{0{,}35,0{,}65\}$, $\rho\in\{0{,}70,0{,}85\}$ y diez semillas por
celda. La factorización empleó una variable por robot y un factor global de
cuota; el pseudoárbol siguió ese orden determinista. En esa factorización, el
factor global induce ancho $N-1$, de modo que el tamaño de UTIL depende del
pseudoárbol empleado.

DPOP coincidió con la factibilidad del MILP en
\NFourDpopMatches{}/\NFourDpopWorlds{} mundos y recuperó el óptimo entero en las
instancias enumeradas. El error objetivo máximo fue $\NFourDpopMaxError{}$ y la
tabla mayor tuvo \NFourDpopMaxEntries{} entradas. La memoria crece
exponencialmente con ese ancho; el experimento acredita exactitud en pequeño, no
escalabilidad.

\clearpage

% --------------------------------------------------------------------------
% Página 38 - escala, tabla maestra, conclusión y SP2
% --------------------------------------------------------------------------
\leveltag{RESULTADOS SP1 · N4 · E6 · SÍNTESIS Y TRANSICIÓN}
\section*{E6: coste creciente en el rango medido}
\enlargethispage{8\baselineskip}

El barrido usa $(N,K)\in\{(16,4),(32,8),(48,12),(64,16)\}$, las geometrías
aleatoria, agrupada y pasillo, cinco semillas por celda, $CV=0{,}65$,
$\rho=0{,}80$ y el grafo medio de N3. Cada método dispone de 4.096 activaciones
como máximo. Son 60 mundos pareados ejecutados en el mismo entorno descrito para
el resto del capítulo. Se registran por separado el tiempo del proceso y los
bytes modelados por agente; el ajuste log--log solo describe ese intervalo.

\begin{figure}[H]
\centering
\includegraphics[width=.60\linewidth]{assets/figures/n4-v2/n4_scaling.pdf}
\caption{E6: escalado observado de BR y DMIS+TX. Los ajustes log--log describen
$16\leq N\leq64$; no son cotas asintóticas.}
\label{fig:n4-scaling-results}
\viuownsource
\end{figure}

Para DMIS+TX, la pendiente temporal observada fue \NFourScaleSlopeD{} (IC
95\,\%: \NFourScaleSlopeDLow{} a \NFourScaleSlopeDHigh{};
$R^2=\NFourScaleRSquaredD$), frente a \NFourScaleSlopeU{} para BR
(IC 95\,\%: 1,59 a 2,19; $R^2=0{,}997$). En $N=64$,
DMIS+TX registró \NFourScaleRuntimeDMax{} ms y \NFourScaleBytesDMax{} bytes por
agente; BR, \NFourScaleRuntimeUMax{} ms y \NFourScaleBytesUMax{}. La curva no se
aplanó. En este intervalo no hay evidencia de coste fijo ni base para extrapolar
a una flota ilimitada.

Ampliar $h$ mejoró el mínimo local; DMIS+TX retiró el orden global a cambio de
tráfico; y DPOP fue exacto solo en pequeño. La Tabla~\ref{tab:n4-family-master}
separa las propiedades formales de esos resultados experimentales antes de
comparar F-II y F-III.

\begin{table}[H]
\centering
\caption{Resultados formales de F-I y alcance de cada enunciado.}
\label{tab:n4-family-master}
\fontsize{7.5}{8.1}\selectfont
\setlength{\tabcolsep}{1.8pt}
\renewcommand{\arraystretch}{0.94}
\begin{tabularx}{\textwidth}{>{\raggedright\arraybackslash}p{2.05cm}
 >{\raggedright\arraybackslash}p{1.55cm}
 >{\raggedright\arraybackslash}p{4.35cm}Y}
\toprule
Resultado & Soporte & Afirmación & Alcance \\
\midrule
Cambio unilateral & identidad~\eqref{eq:n4-unilateral-delta} & El cambio unilateral
se valora con agregados de las cargas afectadas. & Registros vigentes. \\
Potencial exacto & Prop.~4 & La utilidad marginal induce un potencial exacto en cada
fase unilateral. & No cubre $h>1$. \\
Barrera bilateral & Lema~6 y E9 & Puede existir una barrera bilateral; E9 midió
el primer escape conectado. & Búsqueda conectada hasta $h=3$. \\
Terminación local & Prop.~5 & Sin truncamiento, las mejoras estrictas terminan
en un mínimo $h$-local. & No implica óptimo global. \\
Composición & Prop.~6 & Las propuestas independientes conservan el descenso de
$\Psi_4$. & Sin robots ni cargas compartidos. \\
Atomicidad & canal nominal & La confirmación versionada escribe todo o nada. &
Solo bajo el canal nominal. \\
\bottomrule
\end{tabularx}
\viuownsource
\end{table}

\clearpage

% --------------------------------------------------------------------------
% Página 39 - F-II: estado continuo frente a cierre
% --------------------------------------------------------------------------
\leveltag{RESULTADOS SP1 · N4 · E7 · F-II POBLACIONAL}
\section*{F-II: estado continuo y coalición después del cierre}

Las cuatro dinámicas recibieron los mismos \NFourVThreeWorlds{} mundos. No son
los 1.200 mundos de N3.E2: E7 reutiliza el banco congelado de N4.E4, generado
con otra corriente de semillas. N3.E2 contiene 1.159 óptimos y 41 infactibles;
este banco contiene 1.160 y 40, de ahí la diferencia entre denominadores. El
estado continuo final se evaluó al cruzar la tolerancia o agotar 3\,000 iteraciones;
después, el mismo $\mathcal R$ produjo la asignación física. No se descartaron
las ejecuciones que alcanzaron el presupuesto.

\begin{figure}[H]
\centering
\includegraphics[width=.98\linewidth]{assets/figures/n4-v3/n4_fii_continuous_atomic.pdf}
\caption{E7--F-II: déficit relajado al terminar y brecha después del cierre. La brecha
se calcula sobre cierres factibles con óptimo MILP certificado; la tasa de
factibilidad se muestra aparte para no ocultar los fallos.}
\label{fig:n4-fii-continuous-atomic}
\viuownsource
\end{figure}

Smith dejó el mayor potencial terminal en 1\,184 mundos y BNN en 16. Ese orden
describe el potencial terminal: Logit cruzó el criterio numérico en el 84,5\,\%,
Replicator en un mundo y Smith/BNN en ninguno. Tras $\mathcal R$, BNN fue el
mejor resultado entero en 622 mundos, Smith en 300, Logit en 234 y Replicator
en 44. El ganador cambió en \NFourFiiMismatchPct\,\% y la correlación de rangos
mediana fue \NFourFiiRankRho.

Entre los 1\,160 mundos con óptimo entero certificado, las cuatro tasas de
factibilidad quedaron entre 96,1\,\% y 96,8\,\%. Las brechas medianas fueron
18,17\,\% para BNN, \NFourSmithGapPct\,\% para Smith, 20,88\,\% para Logit y
21,38\,\% para Replicator. La discrepancia continuo--entero apareció en las
cinco geometrías.

El potencial continuo y la coalición cerrada no ordenaron igual las dinámicas.
Lo observado pertenece al conjunto dinámica--DAC--$\mathcal R$. Como
$\mathcal R$ es heurístico, no se obtiene una garantía de redondeo óptimo.

\clearpage

% --------------------------------------------------------------------------
% Página 40 - F-III: residuos y resultado entero
% --------------------------------------------------------------------------
\leveltag{RESULTADOS SP1 · N4 · E7 · F-III PRIMAL--DUAL}
\section*{F-III: cierre entero y coste de comunicación}

La convergencia se acepta solo cuando los residuos KKT y de consenso caen por
debajo de $5\times10^{-4}$. Un descenso parcial no basta. La
Figura~\ref{fig:n4-fiii-residuals} conserva la traza de una ejecución y las tasas
por celda.

\begin{figure}[H]
\centering
\includegraphics[width=.98\linewidth]{assets/figures/n4-v3/n4_fiii_residuals.pdf}
\caption{E7--F-III: residuos de una ejecución y porcentaje que cruzó
el criterio conjunto según presión y heterogeneidad.}
\label{fig:n4-fiii-residuals}
\viuownsource
\end{figure}

PD-vGNE distribuido cruzó el umbral en \NFourPdConvergencePct\,\% de los 1\,200
mundos; la mediana KKT final fue $\NFourPdKktMedian$. La tasa bajó de 97,3\,\%
con presión 0,70 y $CV=0$ a 86,0\,\% con $CV=1$. Los 97 casos que no cruzaron
el umbral permanecen en el denominador.

Después de $\mathcal R$, el método distribuido alcanzó 99,57\,\% de factibilidad
y brecha mediana \NFourPdGapPct\,\% en el soporte certificado. En presión 0,85,
frente a Smith+$\mathcal R$, la factibilidad aumentó \NFourPdRiskDiffPct{} puntos
(IC 95\,\%: 3,67 a 7,50; Holm $p=1{,}02 \times 10^{-8}$) y la brecha disminuyó
7,39 puntos (IC 95\,\%: 5,07 a 9,65; Holm
$p=2{,}79 \times 10^{-34}$).

La mediana de tráfico fue 2,73 MB modelados por agente, frente a 0,90 MB para
Smith. El mecanismo primal--dual seguido del mismo cierre produjo más
asignaciones factibles que Smith+$\mathcal R$, a cambio de sostener DAC y
consenso dual. Se compararon los dos mecanismos completos; estos datos no aíslan
el efecto de los precios ni prueban convergencia global del predictor--corrector
digital. Estos bytes cubren la dinámica distribuida, pero excluyen el
postproceso central $\mathcal R$; no son un coste extremo a extremo.

Las comparaciones anteriores resuelven una instantánea. E8 añade una secuencia
de llegadas, finalizaciones y fallos para estudiar la decisión recurrente.

\clearpage

% --------------------------------------------------------------------------
% Página 41 - F-IV: campaña temporal
% --------------------------------------------------------------------------
\leveltag{RESULTADOS SP1 · N4 · E8 · F-IV TEMPORAL}
\section*{F-IV: coste descontado y tasa de servicio}

Las \NFourDynamicWorlds{} secuencias se emparejaron por eventos; dentro de cada
mundo solo cambió la política. El oráculo del objetivo escalar conoce el futuro
y marca un techo de información para ese coste. Esa ventaja impide interpretar
su diferencia frente a las políticas
reactivas como una comparación arquitectónica justa.

\begin{figure}[H]
\centering
\includegraphics[width=.98\linewidth]{assets/figures/n4-v3/n4_fiv_dynamic.pdf}
\caption{E8--F-IV: coste descontado por mundo, cambios de coalición y mensajes.
El cero del oráculo del objetivo escalar en mensajes corresponde a su ejecución
central.}
\label{fig:n4-fiv-dynamic}
\viuownsource
\end{figure}

Las medianas fueron \NFourMyopicCost{} para la resolución reactiva,
\NFourWarmCost{} para la inicialización desde la solución anterior y
\NFourOracleCost{} para el oráculo del objetivo escalar. Frente a la
inicialización anterior, su diferencia pareada fue
\NFourOracleWarmDiff{} (IC 95\,\%: $-0{,}486$ a $-0{,}329$; Holm
$p=5{,}59 \times 10^{-9}$). La inicialización anterior empató con la resolución
reactiva en 15 mundos, ganó en
12 y perdió en tres; redujo las medianas de cambios de 12,5 a 11,5 y de mensajes
de 544,5 a 446,5.

El menor coste del oráculo del objetivo escalar no implica dominancia. Su tasa media de servicio fue
95,6\,\%, frente a 97,8\,\% para las dos políticas Geo-QPG: el objetivo escalar
aceptó en algunos episodios demanda incumplida para evitar distancia o cambios.

Conocer los eventos futuros redujo el coste definido, pero no aumentó la tasa de
servicio. El resultado procede de un problema pequeño con eventos exógenos. En
una aplicación operativa, el servicio tendría que entrar como restricción o como
primera prioridad lexicográfica.

\clearpage

% --------------------------------------------------------------------------
% Página 42 - comparación maestra y conclusión
% --------------------------------------------------------------------------
\leveltag{RESULTADOS SP1 · N4 · COMPARACIÓN Y CONCLUSIÓN}
\section*{Síntesis de SP1: del conteo a la coordinación local}
\enlargethispage{1\baselineskip}

Después del cierre atómico, las tres familias estáticas pueden compararse en el
mismo plano de brecha, factibilidad y tráfico
(Figura~\ref{fig:n4-cross-family-pareto}). F-IV queda fuera porque se evaluó con
secuencias independientes.

\begin{figure}[H]
\centering
\includegraphics[width=.58\linewidth]{assets/figures/n4-v3/n4_cross_family_pareto.pdf}
\caption{Síntesis transversal de N3--N4 sobre los mundos comparables: brecha
frente al MILP, factibilidad y tráfico modelado. El tamaño representa
factibilidad; los bytes cubren la fase decisional y excluyen el cierre central
$\mathcal R$, por lo que no son un coste extremo a extremo.}
\label{fig:n4-cross-family-pareto}
\viuownsource
\end{figure}

\begin{table}[H]
\centering
\caption{Síntesis de N1--N4 y razón de cada transición.}
\label{tab:sp1-level-synthesis}
\fontsize{7.1}{7.5}\selectfont
\setlength{\tabcolsep}{1.7pt}
\renewcommand{\arraystretch}{0.88}
\begin{tabularx}{\textwidth}{>{\raggedright\arraybackslash}p{.48cm}
 >{\raggedright\arraybackslash}p{2.00cm}
 >{\raggedright\arraybackslash}p{1.95cm}
 >{\raggedright\arraybackslash}p{3.70cm}Y}
\toprule
 & Pregunta & Método y arquitectura & Resultado que queda establecido & Frontera \\
\midrule
N1 & ¿Cuántos robots? & LSAP central; flota homogénea & Expansión exacta por
puestos. & Falla al variar la capacidad. \\
N2 & ¿Quién cubre la carga? & MILP central; estado global & Coalición entera y
oráculo certificado. & Certificación temporal y coordinador. \\
N3 & ¿Quién decide? & CBBA-RB / GRAPE; vecinos & GRAPE gana factibilidad;
CBBA-RB ahorra tráfico. & Conectividad y coste de reconstruir el perfil. \\
N4 & ¿Qué orden local? & Geo-QPG F-I; vecinos & 1\,156/1\,200 terminales BR
tuvieron primero un escape bilateral. & $h\le3$, canal nominal y mundo estático. \\
\bottomrule
\end{tabularx}
\viuownsource
\end{table}

Los pares no explican el 96,3\,\% del beneficio. Ese porcentaje corresponde a
los 1\,156 terminales BR cuyo primer escape, buscado hasta $h=3$, fue bilateral.
En este régimen tiene sentido probar pares antes que órdenes superiores. Nada de
ello cubre un canal degradado ni una ejecución física.

\begin{tcolorbox}[colback=SPFloor,colframe=VIUOrange,boxrule=.7pt,arc=1mm,
 left=2.3mm,right=2.3mm,top=1.2mm,bottom=1.2mm]
\footnotesize
\textbf{Conclusión.} SP1 entrega una coalición entera cuya cobertura se verifica
antes de SP2; las demás ejecuciones se registran como no factibles. La secuencia
N1--N4 abandona el recuento homogéneo, incorpora capacidad individual y retira
el coordinador. En N4, coordinar dos o tres robots mejoró mínimos que no cedían
al cambiar solo la regla unilateral. Geo-QPG F-I concentra la aportación; F-II y
F-III son contrastes continuos bajo un cierre común, y F-IV delimita la extensión
temporal. Acoplamiento, \emph{wrench} y transporte corresponden a SP2.
\end{tcolorbox}
